\def\fileversion{v1.0}
\def\filename{livroaberto-professor}
\def\filetitle{Livro Aberto package}
\def\filedate{2018/02/13}
\def\docdate {2018/02/13}

%Combinação dos pacotes usados para o livro aberto
%Para acessar cada seção basta usar Ctrl+F

%Seção 0 - Pacotes e Comandos base
%Seção 1 - Livro Aberto Boxes
%Seção 2 - Livro Aberto Colors
%Seção 3 - Livro Aberto Fonts
%Seção 4 - Livro Aberto Headings
%Seção 5 - Livro Aberto Lists
%Seção 6 - Livro Aberto Page 
%Seção 7 - Livro Aberto Professor


% Define o diretório das figuras e das fontes
\def\fontspath{Fontes/}
\def\graphics{Figuras/}

\def\aluno{0}


%Base Packages
\RequirePackage[skins,breakable]{tcolorbox}
\RequirePackage{varwidth}
\RequirePackage{array}
\RequirePackage{colortbl}
\RequirePackage{bookmark}
\RequirePackage{hyperref}
\PassOptionsToPackage{}{hyperref}
\hypersetup{colorlinks=true, colorlinks=true, linkcolor=linkscolor, filecolor=linkscolor, urlcolor=linkscolor,citecolor=linkscolor}
\RequirePackage{float}
\RequirePackage{tkz-euclide}
\RequirePackage{amsmath,amssymb}
\RequirePackage[brazil]{babel}
\RequirePackage{times}
\RequirePackage{tabulary}
\RequirePackage{geometry}
\RequirePackage{hypcap}
\RequirePackage{makeidx}
\RequirePackage{xurl}
\RequirePackage{amsmath,amssymb,amsfonts,amsthm,amstext}
\RequirePackage{tabu} 
\RequirePackage{hhline}
\RequirePackage{pbox}
\RequirePackage{ragged2e}
\RequirePackage{microtype}
\RequirePackage{tabularx}
\RequirePackage{makecell}
\RequirePackage{longtable}
\RequirePackage{wrapfig}
\RequirePackage{bm}
\RequirePackage[low-sup]{subdepth} 
\RequirePackage{tikz-3dplot}
\RequirePackage{pgfplots}
\RequirePackage[inline, attach]{asymptote}
\RequirePackage[dvips]{attachfile2}
\RequirePackage{vwcol}
\RequirePackage{tikz}
\RequirePackage{ifthen}
\RequirePackage[export]{adjustbox}
\RequirePackage{xparse}% For RenewDocumentEnvironment
\RequirePackage{expl3}
\RequirePackage{tcolorbox} %Allow those nice environments
\tcbuselibrary{theorems, listings, breakable, most, skins}
\RequirePackage{titling} % needed to manage with the title
\RequirePackage{eso-pic }% to the cover style
\RequirePackage{tikz}
\usetikzlibrary{intersections, positioning, fit, backgrounds, shapes, fadings,
decorations.pathmorphing, graphs, quotes, angles, calc, through, backgrounds}
\RequirePackage{comment}
\RequirePackage{ltxcmds}
\RequirePackage{amstext}
\RequirePackage{textcomp}% "warn" option issued from template
% \RequirePackage[explicit]{titlesec}
\PassOptionsToPackage{explicit}{titlesec}
\RequirePackage{titletoc}
\RequirePackage{longtable}
\RequirePackage{makeidx}
\RequirePackage{fancyvrb}
\RequirePackage{wrapfig}
\RequirePackage{parskip}[=v1]
\RequirePackage{parskip}
\RequirePackage{alltt}
\RequirePackage{upquote}
\RequirePackage{capt-of}
\RequirePackage{needspace}
\RequirePackage{mathtools}
\ltx@ifundefined{@removefromreset}
{\RequirePackage{remreset}}
\RequirePackage{atbegshi}
\RequirePackage{kvoptions}
\RequirePackage{sphinx}
\@ifclassloaded{memoir}{}{\RequirePackage{fancyhdr}}
\RequirePackage{etoolbox}
\RequirePackage{framed}
\RequirePackage{siunitx}
\RequirePackage{caption}
\RequirePackage{pdfpages}
\RequirePackage{layout}
\RequirePackage{stackengine}
\RequirePackage{changepage}
\RequirePackage{graphicx}
\RequirePackage{adjustbox}
\RequirePackage{titlecaps}
% \RequirePackage{scrlayer-notecolumn}
\RequirePackage{marginfix}
\RequirePackage{pdfpages}
\RequirePackage{stmaryrd}
\RequirePackage{lscape}

\RequirePackage[sectionbib]{natbib}
\RequirePackage{chapterbib}
\patchcmd{\thebibliography}{%
  \chapter*{\bibname}\@mkboth{\MakeUppercase\bibname}{\MakeUppercase\bibname}}{%
  \clearpage
  \bibliographytitle*{Bibliografia}}{}{}

\RequirePackage{spverbatim}
% \RequirePackage{ulem}
\RequirePackage{cancel}
\RequirePackage{ifthen}
\RequirePackage{booktabs}
\RequirePackage{listofitems}
\RequirePackage{url}
\RequirePackage{graphicx}
%\RequirePackage[usenames,dvipsnames]{color}
\RequirePackage{forloop}
\RequirePackage{lipsum}
\RequirePackage{suffix}
\RequirePackage{refcount}
% \RequirePackage[hidelinks]{hypreref}
\RequirePackage{xcolor}
\RequirePackage{multicol}
\RequirePackage{etex}
\RequirePackage{multirow}
%\RequirePackage[brazilian]{babel}
\RequirePackage{blindtext}

\RequirePackage{cleveref}
\RequirePackage{collcell}
\RequirePackage{datatool}
\RequirePackage{environ}
\RequirePackage{xstring}

\tcbuselibrary{skins,breakable}
\RequirePackage{graphicx}
\@ifclassloaded{memoir}{}{\RequirePackage{fancyhdr}}

\DeclareMathOperator{\sen}{sen}
\DeclareMathOperator{\tg}{tg}
\DeclareMathOperator{\cotg}{cotg}
\DeclareMathOperator{\cosec}{cosec}
\DeclareMathOperator{\arcsen}{arcsen}
\DeclareMathOperator{\arctg}{arctg}
\DeclareMathOperator{\senh}{senh}
\DeclareMathOperator{\tgh}{tgh}
\DeclareMathOperator{\mdc}{mdc}
\DeclareMathOperator{\mmc}{mmc}
% \DeclareMathOperator{\si}{Si}
\DeclareMathOperator{\dirac}{\delta}
\DeclareMathOperator{\mmu}{\mu}
\DeclareMathOperator{\rad}{rad}

\newcommand{\R}{\mathbb{R}}
\newcommand{\Q}{\mathbb{Q}}
\newcommand{\Z}{\mathbb{Z}}
\newcommand{\N}{\mathbb{N}}

%Commands
\renewcommand{\bottomtitlespace}{0.25\textheight}
%\pgfplotsset{compat=1.15}
\renewcommand{\thead}{\rowfont{\bfseries \color{white}} \rowcolor{\currentcolor!80}}
\newcommand{\tnumber}{\rowfont{\color{white}}\rowcolor{\currentcolor!80}}

\newcommand{\tmat}[1]{\cellcolor{\currentcolor!80}\textcolor{white}{\bm{#1}}}
\newcommand{\tcolor}[1]{\cellcolor{\currentcolor!80}\textcolor{white}{\textbf{#1}}}
\newcommand{\tmcol}[3]{\multicolumn{#1}{#2}{\cellcolor{\currentcolor!80}\textcolor{white}{\textbf{#3}}}}
\newcommand{\titem}[1]{\textcolor{\currentcolor}{\textbf{#1}}}

\newcommand{\tikzcolor}{\currentcolor!80}
\newcommand{\sub}{\textsubscript}
\newcommand{\super}{\textsuperscript}

\newcommand{\newparallel}{\mathrel{\mkern-5mu\clipbox{0 0.75ex 0 0}{${\sslash}$}\!}}

\newcolumntype{d}[1]{>{\vspace{2.5pt}}m{#1}<{\vspace{2.5pt}}}
\newcolumntype{e}[1]{>{\vspace{2.5pt}\centering}m{#1}<{\vspace{2.5pt}}}
\newcolumntype{f}{>$c<$}


%Coloca listas de enumerate no padrão da seção de exercícios
\newcommand\exerciselist
{	
	\setlist[enumerate]{label=\protect\enumbox{\arabic*}}
	\setlist[enumerate,2]{label=\textit{\alph*)}}
	\setlist[enumerate,3]{label=\textit{\roman*)}}
}





%Seção 1 - Livro Aberto Boxes
%***************************************************************************************
% This package creates and typesets the text boxes for the Livro Aberto textbook.
%
% This package was written and created by
% David Latchman, TeXnical Designs
% Contact: david.latchman@texnical-designs.com for problems

% \def\fileversion{v1.0}
% \def\filename{livroabertoboxes}
% \def\filetitle{Livro Aberto Text Boxes package}
% \def\filedate{2018/02/13}
% \def\docdate {2018/02/13}


\NeedsTeXFormat{LaTeX2e}
\ProvidesPackage{\filename}

\typeout{Package: `\filetitle\space\filename\space\fileversion \@spaces <\filedate>'}
\typeout{English documentation as of <\docdate>}

%% ---------------------------------------------------------------------


%%
%% Loads the ifthen package
%%
\RequirePackage{ifthen}
\RequirePackage{mfirstuc} %To capitalize just the inicial letter of each word. There is a command to exclude specific words to captalize such like a, an, the, etc.
\RequirePackage{anyfontsize} %To define fontsize precisely

\tikzstyle{outlinebox} = [fill=box2]
\tikzstyle{box} = [inner sep=3mm, inner ysep=3mm]

% Define the vertical skip for the environments
\newskip\libroskip
\libroskip= 4pt plus 1pt minus 1pt
\hfuzz=4pt
\graphicspath{\graphics}


\def\envskip{\vskip\libroskip}



\newtcolorbox[auto counter, number within=chapter, number freestyle={\noexpand\arabic{\tcbcounter}}]{task}[1]{parbox=false, blanker, enhanced, breakable, before skip = 8mm, after skip = 5mm, left=0.3mm, right=0.8mm, top=5mm, bottom=3mm, colback=white, colframe=cor2!70, width=162mm, 
%borderline west={3mm}{0pt}{white}, 
  toprule=0pt,
  bottomrule=0.5pt,
  rightrule=-1pt,
  leftrule=-1pt,
  outer arc = 3mm,
  arc = 3mm,
  sharp corners = northeast,
  sharp corners = southeast,
  sharp corners = southwest,
  fontupper=\fontsize{11}{15.5}\boxexample, % Should be 11 pt.
fonttitle=\boxtitlefont\large,  % Should be 9 pt. Don't know how to set this. Tried \fontsize{9pt}{}, but didn't work.
coltitle=\currentcolor,
title=\makefirstuc{#1\hfill}, 
boxed title style={empty,arc=1pt,outer arc=1pt,boxrule=1pt}, 
attach boxed title to top left={xshift=-1 mm , yshift=-0.1mm},
underlay boxed title ={
% Draw the small arc on the top left.
\draw[cor2!70, line width = 1pt] ($(title.south west)+(3mm,0mm)$) coordinate (LU) arc [start angle=90, end angle=180, y radius=2mm, x radius=2mm] --++ (0mm , -2mm);
% Draw the top line from the left to the right.
\draw[cor2!70,  line width=1pt] (LU) -- ++ (158mm,0mm) coordinate (R);
% Think on the Atividade box as an ABCD rectangle (clockwise, A is top left), now you see the coordinates. 
\fill [cor2!70] ($(R)+(0mm,-.3mm)$) -- ++ (-25.00mm,0mm) coordinate (D) --++ (0mm, 5mm) arc [start angle=180, end angle=90, x radius=2.00mm, y radius=2.00mm] -- ++ (25.00mm,0mm) coordinate (B) --++ (0.00mm,-5mm)  arc [start angle=0, end angle=-90, y radius=2.00mm, x radius=2.00mm] -- cycle;
\node[font=\titlefont\large,white] at ($(B)!0.5!(D)$) {\small Atividade \thetcbcounter};
}
}


% Makes the "To think about" box.
% It is missing to set title font and font size. After that should check if the title is centralized in the title box.
\newtcolorbox{reflection}{parbox=false, blanker, enhanced, breakable, before skip = 5mm, after skip = 5mm, left=3mm, right=3mm, top=7mm, bottom=3mm, colback=white, colframe=box1, borderline west={3mm}{0pt}{white}, toprule=1pt, bottomrule=1pt, rightrule=1pt, leftrule=-1pt, outer arc = 3mm, arc = 3mm, 
fonttitle=\bfseries\large\boxtitlefont,  % Should be 12 pt. Don't know how to set this.
title=PARA REFLETIR, 
boxed title style={empty,arc=0pt,outer arc=0pt,boxrule=0pt}, 
attach boxed title to top left={xshift=3mm , yshift=-2.5mm},
underlay boxed title={
\filldraw [draw=box1, fill=box1] ($(title.north west) + (0 mm,-2.50mm)$) -- ++(0.00mm,1.50mm) arc [start angle=180, end angle=90, y radius=2.00mm, x radius=2.00mm] -| ++(34.00mm,-5.00mm) arc [start angle=0, end angle=-90, x radius=2.00mm, y radius=2.00mm] -| cycle;}
}



% Observation Box

\newtcolorbox{observationtitle}[1]{
parbox=false, enhanced, breakable, left=3mm, right=3mm,bottom=2mm, colback=box2!50, colframe=box2!50, arc=3mm, sharp corners=uphill, pad at break=2mm, drop fuzzy shadow, title={\bfseries\textcolor{\currentcolor}{ #1}}, boxed title style ={empty},top=8mm, attach boxed title to top center={xshift=0mm, yshift=-7.5mm}
}

\newtcolorbox{observation}[1]{parbox=false, enhanced, breakable, left=3mm, right=3mm, top=3mm, bottom=2mm, colback=box2!50, colframe=box2!50, arc=3mm, sharp corners=uphill, pad at break=2mm, drop fuzzy shadow, title={\bfseries\textcolor{\currentcolor}{ #1}}, attach boxed title to top center={xshift=0mm, yshift=-5.5mm}, boxed title style ={empty}}	 

\newtcolorbox[auto counter, number within=chapter, number freestyle={\noexpand\arabic{\tcbcounter}}]{example}[1]{parbox=false, blanker, enhanced, breakable, before skip = 5mm, after skip = 5mm, left=3mm, right=3mm, top=10mm, bottom=3mm, colback=white, colframe=cor1, width=162mm, toprule=1pt, bottomrule=1pt, rightrule=1pt, leftrule=1pt, outer arc = 8mm, arc = 8mm, sharp corners = northwest, sharp corners = southeast, sharp corners = southwest,
title={
%{\boxfont EXEMPLO}\hspace{3.5mm} 
\textcolor{cor2}{\large\extitlefont\makefirstuc{#1}\phantom{ç}}
}, 
boxed title style={empty, sharp corners}, 
attach boxed title to top left={xshift=22mm , yshift=-7.5mm},
underlay boxed title ={
% Draw the title underlay rectangle
\def\temp{#1}\ifx\temp\empty
%
\else
\fill[cor1!30] ($(frame.north west)+(0,-1pt)$) rectangle ($(title.south east)+(3mm,0mm)$) coordinate (R);
\fi
%)
%;
% Filldraw the darker fancy in front of the rectangle and behind the Example. The coordinates should be better chosen.


\fill[cor1] (frame.north west) -- ($(frame.north west)+(0,-1.4pt)+(0,-7mm)$) -- ++ (22mm,0mm) arc [start angle=-90, end angle=0, x radius = 2mm, y radius = 2mm] -- ($(frame.north west) + (24mm,0)$) --  cycle;
\node [white] at ($(frame.north west) + (12mm,-3.5mm)$) {\boxfont EXEMPLO \thetcbcounter};

},
before upper={\def\currentcolor{cor1}}
}

\newtcolorbox{research}{parbox=false, blanker, enhanced, breakable, before skip = 5mm, after skip = 5mm, left=3mm, right=3mm, top=6mm, bottom=3mm, colback=white, colframe=box3, toprule=1pt, bottomrule=1pt, rightrule=1pt, leftrule=1pt, outer arc = 3mm, arc = 3mm, sharp corners = northwest,
fonttitle=\normalsize\boxtitlefont,  
title= PARA PESQUISAR, 
boxed title style={empty, colback = box3, colframe = box3, sharp corners}, 
attach boxed title to top left={xshift=-2.5mm , yshift=-6.2mm},
underlay boxed title={
\filldraw [fill=box3, draw=box3] ($(frame.north west)$) rectangle +(36.00mm,-5.00mm);% the rectangle where the title fits.
\coordinate (A) at ($(frame.north west) + (36.00mm,0.00mm)$);
%\filldraw [cor1] ($(A) + (-4.7mm,0mm)$) rectangle ++(-7mm,-6.4mm);%green rectangle within ? inside.
%\fill [box3] ($(A) + (-11.9mm,-3.5mm)$) -- ++(4.8 mm,2.6 mm) arc [start angle=20, end angle=-60, x radius=1.00mm, y radius=2.35mm] --++ (0,-2mm)-- ($(title.south east) + (-1.1mm,0.20mm)$) -- ($(title.south east) + (-2mm,0.20mm)$) -- cycle; % the fancy thing in the ?
\node [white] at ($(A) + (-2.40mm,-1.6mm)$) {\includegraphics[height=7mm]{lupa}};% just the ?
}
}



\newtcolorbox{knowledge}{parbox=false, blanker, enhanced, breakable, before skip = 5mm, after skip = 5mm, left=3mm, right=3mm, top=6mm, bottom=3mm, colback=white, colframe=box3, toprule=1pt, bottomrule=1pt, rightrule=1pt, leftrule=1pt, outer arc = 3mm, arc = 3mm, sharp corners = northwest,
fonttitle=\normalsize\boxtitlefont,  
title= VOC\^{E} SABIA, 
boxed title style={empty, colback = box3, colframe = box3, sharp corners}, 
attach boxed title to top left={xshift=-2.3mm , yshift=-5.5mm},
underlay boxed title={
\filldraw [fill=box3, draw=box3] ($(frame.north west) + (.4mm,0mm)$) rectangle ++(24.00mm,-5.00mm);% the rectangle where the title fits.
\coordinate (A) at ($(frame.north west) + (24.00mm,-2mm)$);
\node at (A) {\includegraphics[scale=1]{interrogacao.png}};
}
}


\NewEnviron{evaluation}{%
    \envskip
	\begin{tikzpicture}[remember picture]
		\begin{pgfonlayer}{main}
			\node [box] (box) {%
			\begin{minipage}{0.90\textwidth}
				\BODY			
			\end{minipage}	
			};
		\end{pgfonlayer}
		\begin{pgfonlayer}{main}
			\filldraw [fill=cor1, draw=cor1] (box.north west) arc [start angle=270, end angle=180, x radius=2.00mm, y radius=2.00mm] -- ++(0.00mm,6.00mm) arc [start angle=180, end angle=90, x radius=2.00mm, y radius=2.00mm] -- ++(3.00mm,0.00mm) coordinate (a) |- cycle;
			\node [draw=cor1, fill=cor1, minimum height=10.00mm, right] (heading) at ($(a) + (0.00mm,-5.00mm)$) {\Large AUTO-AVALIA\c{C}\={A}O};
			\filldraw [fill=cor1, draw=cor1] (heading.north east) -- ++(10.00mm,-10.00mm) coordinate (b) -| cycle;
		\end{pgfonlayer}
		
		\begin{pgfonlayer}{background}
			\draw [draw=cor1!30, line width=3.00mm, overlay] ($(box.north west) + (0.00mm,1.50mm)$) -| ($(box.south east) + (1.50mm,-1.50mm)$) -- ($(box.south west) + (-1.50mm,-1.50mm)$);
			\draw [draw=cor1, overlay] (b) -- (box.north east) -- (box.south east) -- ($(box.south west) + (-1.50mm,0.00mm)$);

%			\draw [draw=cor1!30, line width=3.00mm, overlay] ($(box.north west) + (0.00mm,1.50mm)$) -| ($(box.south east) + (1.50mm,-1.50mm)$) coordinate (d) -- ($(current page.north west)!(d)!(current page.south west)$) -- (d);
%			\draw [draw=cor1, overlay] (b) -- (box.north east) -- (box.south east) -- ($(box.south west) + (-15.00mm,0.00mm)$) coordinate (c) -- ($(current page.north west)!(c)!(current page.south west)$) -- (c);
		\end{pgfonlayer}	
		
	\end{tikzpicture}

}

\NewEnviron{project}{%
\def\currentcolor{cor2}
    \newpage
	\begin{tikzpicture}[remember picture, overlay]
		\coordinate (border1) at ($(current page text area.south east) + (0.00mm,8.00mm)$);
		\coordinate (border1a) at ($(border1) + (-5.00mm,2.00mm)$);
		\coordinate (border1b) at ($(current page text area.south west) + (11.00mm,10.00mm)$);
		\coordinate (border1c) at ($(current page text area.north west) + (11.00mm,-13.00mm)$);
        \coordinate (border1d) at ($(current page text area.south west) + (11.00mm,13.00mm)$);
		\coordinate (heading) at ($(current page text area.north west) + (32.50mm,-5.00mm)$);
		\coordinate (corner1) at ($(current page text area.north west) + (6.00mm,-10.00mm)$);
		\coordinate (corner2) at ($(current page text area.north east) + (-4.00mm,-10.00mm)$);
		\coordinate (corner3) at ($(current page text area.south east) + (0.00mm,10.00mm)$);
		\coordinate (corner4) at ($(current page text area.south west) + (10.00mm,6.00mm)$);
		\coordinate (text1) at ($(current page text area.north west) + (11.00,-12.50mm)$);
		\coordinate (text2) at ($(current page text area.south east) + (-2.50mm,8.50mm)$);
        \coordinate (bodytextpos) at ($(border1c) + (5.00mm,-3.00mm)$);
		
		\begin{pgfonlayer}{background}
			\draw [draw=base1, fill=base1] (current page text area.south west) rectangle (current page text area.north east);
            		\end{pgfonlayer}
		
		\begin{pgfonlayer}{main}
			\draw [fill=cor2, draw=cor2] (current page text area.north west) -- ++(0.00mm,-8.00mm) arc [start angle=180, end angle=270, x radius=2.00mm, y radius=2.00mm] -| ++(63.00mm,8.00mm) arc [start angle=0, end angle=90, x radius=2.00mm, y radius=2.00mm] -- cycle;
			\draw [fill=base1!30, draw=base1!30] (corner1) -- (corner2) arc [start angle=90, end angle=0, x radius=4.00mm, y radius=4.00mm] -- (corner3) arc [start angle=0, end angle=-90, x radius=4.00mm, y radius=4.00mm] -- (corner4) arc [start angle=270, end angle=180, x radius=4.00mm, y radius=4.00mm] -- cycle;
			% \draw [green] (text1) rectangle (text2);
			\draw [black] (border1a) -- (border1b) -- (border1c) -| cycle;

		\end{pgfonlayer}
		
		\begin{pgfonlayer}{foreground}
			\node [color=base1!30, font=\titlefont\bfseries\fontsize{13pt}{0}\selectfont] (title) at (heading) {PROJETO APLICADO};
            \node [box, below right] (bodytext) at (border1c){%
			\begin{minipage}{0.85\textwidth}
				\BODY			
			\end{minipage}	
			};
		\end{pgfonlayer}
	
		
	\end{tikzpicture}
    \clearpage
}

%Seção 2 - Livro Aberto Colors

%***************************************************************************************
% This package sets up the colors used in the Livro Aberto textbook.
%
% This package was written and created by
% David Latchman, TeXnical Designs
% Contact: david.latchman@texnical-designs.com for problems
%
%
% The fonts used in this package:
% Title: Orbitron
% Main (Text): Exo 2
% Boxes: Roboto

% \def\fileversion{v1.0}
% \def\filename{livroabertocolors}
% \def\filetitle{Livro Aberto Colors package}
% \def\filedate{2018/02/13}
% \def\docdate {2018/02/13}


\NeedsTeXFormat{LaTeX2e}
\ProvidesPackage{\filename}

\typeout{Package: `\filetitle\space\filename\space\fileversion \@spaces <\filedate>'}
\typeout{English documentation as of <\docdate>}

%% ---------------------------------------------------------------------
%%
%% Loads the ifthen package
%%
\RequirePackage{ifthen}

\RequirePackage{xkeyval}

\definecolor{cor1}{RGB}{0,165,157} %verde agua
\definecolor{cor2}{RGB}{73,85,93}
\definecolor{cor1a}{RGB}{0,165,157}
\definecolor{cor2a}{RGB}{73,85,93}
\definecolor{cor1b}{RGB}{251,181,0}
\definecolor{cor2b}{RGB}{73,85,93}

\definecolor{base1}{RGB}{184,202,213}
\definecolor{base2}{RGB}{221,223,80}

\definecolor{session1}{RGB}{0,94,176} %azul
\definecolor{session2}{RGB}{75,173,49} %verde
\definecolor{session3}{RGB}{183,13,40} %vinho
\definecolor{session4}{RGB}{236,107,16} %laranja

\definecolor{penaqua}{RGB}{0,165,157}
\definecolor{grafite}{RGB}{73,85,93}
\definecolor{penazul}{RGB}{0,94,176}
\definecolor{penverde}{RGB}{75,173,49}
\definecolor{penvinho}{RGB}{83,13,40}
\definecolor{penlaranja}{RGB}{236,107,16}
\definecolor{penbox2}{RGB}{240,238,201}

\definecolor{box1}{RGB}{71,30,129}
\definecolor{box2}{RGB}{225,221,146}
\definecolor{box3}{RGB}{108,27,98}
\definecolor{linkscolor}{RGB}{6,69,173}

\definecolor{magblue}{RGB}{151,255,249}

\definecolor{primario}{RGB}{0,165,157} %esverdeado
\definecolor{secundario}{RGB}{73,85,93} %acinzentado escuro
\definecolor{terciario}{RGB}{184,202,213} % azul claro
\definecolor{destacado}{RGB}{183,13,40} %avermelhado
\definecolor{atento}{RGB}{0,94,176} %azul forte

% Set up options for color schemes
\newboolean{coloroptions}

\define@choicekey*+{colorchoices}{color options}[\val\coloroption]
    {option 1, option 2}{%
            \ifcase\coloroption\relax
                \definecolor{cor1}{RGB}{0,165,157}
                \definecolor{cor2}{RGB}{73,85,93}
                % \setboolean{coloroptions}{true}
            \or
                % \setboolean{coloroptions}{false}
                \definecolor{cor1}{RGB}{251,181,0}
                \definecolor{cor2}{RGB}{73,85,93}
            \fi
        }{%
            % \setboolean{coloroptions}{true}
            \definecolor{cor1}{RGB}{0,165,157}
            \definecolor{cor2}{RGB}{73,85,93}
        }

%Seção 3 - Livro Aberto Fonts

%***************************************************************************************
% This package sets up the fonts for the Livro Aberto textbook.
%
% This package was written and created by
% David Latchman, TeXnical Designs
% Contact: david.latchman@texnical-designs.com for problems
%
%
% The fonts used in this package:
% Title: Orbitron
% Main (Text): Exo 2
% Boxes: Roboto

% \def\fileversion{v1.0}
% \def\filename{livroabertofonts}
% \def\filetitle{Livro Aberto Fonts package}
% \def\filedate{2018/02/13}
% \def\docdate {2018/02/13}


\NeedsTeXFormat{LaTeX2e}
\ProvidesPackage{\filename}

\typeout{Package: `\filetitle\space\filename\space\fileversion \@spaces <\filedate>'}
\typeout{English documentation as of <\docdate>}

%% ---------------------------------------------------------------------
%%
%% Loads the ifthen package
%%
\RequirePackage{ifthen}

\def\fonts{Fontes/}

\RequirePackage{fontspec}
\setmainfont[Path=\fonts]{Exo2-Light.otf}[BoldFont=Exo2-Medium.otf, ItalicFont=Exo2-LightItalic.ttf, BoldItalicFont=Exo2-MediumItalic.ttf]
\newfontfamily{\titlefont}[Path=\fonts]{orbitron-medium.otf}[BoldFont=orbitron-bold.otf]
\newfontfamily{\secondtitlefont}[Path=\fonts]{orbitron-black.ttf}
\newfontfamily{\thirdtitlefont}[Path=\fonts]{orbitron-bold.otf}
\newfontfamily{\pagetitlefont}[Path=\fonts]{Orbitron-Regular.ttf}
\newfontfamily{\boxfont}[Path=\fonts]{Roboto-Regular.ttf}
\newfontfamily{\enumeratefont}[Path=\fonts]{Exo2-Light.otf}[BoldFont=Exo2-Medium.otf, ItalicFont=Exo2-LightItalic.ttf, BoldItalicFont=Exo2-MediumItalic.ttf]
\newfontfamily{\boxtitlefont}[Path=\fonts]{Exo2-Bold.otf}[BoldFont=Exo2-ExtraBold.ttf, ItalicFont=Exo2-BoldItalic.ttf, BoldItalicFont=Exo2-ExtraBoldItalic.ttf] % added after David
\newfontfamily{\extitlefont}[Path=\fonts]{Exo2-SemiBold.otf} % added after David
\newfontfamily{\boxexample}[Path=\fonts]{Exo2-Regular.otf}[BoldFont=Exo2-SemiBold.ttf, ItalicFont=Exo2-RegularItalic.ttf, BoldItalicFont=Exo2-MediumItalic.ttf] % added after David

%Seção 4 - Livro Aberto Headings

%***************************************************************************************
% This package creates the heading designs for the Livro Aberto textbook.
%
% This package was written and created by
% David Latchman, TeXnical Designs
% Contact: david.latchman@texnical-designs.com for problems
%
%

% \def\fileversion{v1.0}
% \def\filename{livroabertoheadings}
% \def\filetitle{Livro Aberto Headings package}
% \def\filedate{2018/02/13}
% \def\docdate {2018/02/13}


\NeedsTeXFormat{LaTeX2e}
\ProvidesPackage{\filename}

\typeout{Package: `\filetitle\space\filename\space\fileversion \@spaces <\filedate>'}
\typeout{English documentation as of <\docdate>}

%% ---------------------------------------------------------------------
%%
%% Loads the ifthen package
%%
\RequirePackage{ifthen}

% Commands for chapter page
\newcommand\chapterillustration{}
\newcommand\chapterwhat{}
\newcommand\chapterbecause{}

% Formatação dos comandos \section e \subsection
%Cor atual
\titleformat{\section}{\color{\currentcolor}\normalfont\semihuge\bfseries}{\thesection}{1em}{\phantomsection #1}
\titlespacing*{\section}
{0pt}{2mm}{1mm}
% \providecommand*{\toclevel@section}{2}


\titleformat{\subsection}{\color{\currentcolor}\normalfont\Large\bfseries}{\thesection}{1em}{#1}
\titlespacing*{\subsection}
{0pt}{2mm}{1mm}

\titleformat{\paragraph}{\vspace{.3\baselineskip}\color{\currentcolor}\normalfont\large\bfseries}{\thesubsection}{1em}{#1}
\titlespacing*{\paragraph}
{0pt}{1mm}{1mm}



\renewcommand*{\toclevel@section}{1}
\renewcommand*{\toclevel@subsection}{4}
\renewcommand*{\toclevel@paragraph}{5}
\renewcommand*{\toclevel@subparagraph}{6}

\providecommand*{\toclevel@exploresec}{2}
\providecommand*{\toclevel@practicesec}{2}
\providecommand*{\toclevel@arrangesec}{2}
\providecommand*{\toclevel@knowsec}{1}
\providecommand*{\toclevel@exercisesec}{1}

\titleformat{\bibliographytitle}{\vspace{.3\baselineskip}\color{\currentcolor}\normalfont\huge\bfseries}{\thesubsection}{1em}{}
\titlespacing*{\subsection}
{0pt}{1mm}{1mm}

%Cor Explorando

\newcommand{\scexplore}[1]{\hypersetup{linkcolor=session1}\paragraph{\color{session1}\normalfont\Large\bfseries#1}\hypersetup{linkcolor=linkscolor}}
\newcommand{\scarrange}[1]{\hypersetup{linkcolor=session4}\paragraph{\color{session4}\normalfont\Large\bfseries#1}\hypersetup{linkcolor=linkscolor}}
\newcommand{\scpractice}[1]{\hypersetup{linkcolor=session2}\paragraph{\color{session2}\normalfont\Large\bfseries#1}\hypersetup{linkcolor=linkscolor}}
\newcommand{\scknow}[1]{\hypersetup{linkcolor=session2}\paragraph{\color{session3}\normalfont\Large\bfseries#1}\hypersetup{linkcolor=linkscolor}}
\newcommand{\scexercise}[1]{\hypersetup{linkcolor=cor1}\paragraph{\color{cor1}\normalfont\Large\bfseries#1}\hypersetup{linkcolor=linkscolor}}


\newcommand{\subexplore}[1]{\subparagraph{\color{session1}\normalfont\Large\bfseries#1}}
\newcommand{\subarrange}[1]{\subparagraph{\color{session4}\normalfont\Large\bfseries#1}}
\newcommand{\subpractice}[1]{\subparagraph{\color{session2}\normalfont\Large\bfseries#1}}
\newcommand{\subknow}[1]{\subparagraph{\color{session3}\normalfont\Large\bfseries#1}}
\newcommand{\subexercise}[1]{\subparagraph{\color{cor1}\normalfont\Large\bfseries#1}}


% Create box shape
\pgfdeclareshape{howbecausetextbox}{%
    \inheritsavedanchors[from=rectangle]
    \inheritanchorborder[from=rectangle]

    \inheritanchor[from=rectangle]{center}
    \inheritanchor[from=rectangle]{south}
    \inheritanchor[from=rectangle]{west}
    \inheritanchor[from=rectangle]{north}
    \inheritanchor[from=rectangle]{east}
    \inheritanchor[from=rectangle]{base}
    \inheritanchor[from=rectangle]{north west}
    \inheritanchor[from=rectangle]{north east}
    \inheritanchor[from=rectangle]{south east}
    \inheritanchor[from=rectangle]{south west}

    \backgroundpath{%
        \southwest \pgf@xa=\pgf@x \pgf@ya=\pgf@y
        \northeast \pgf@xb=\pgf@x \pgf@yb=\pgf@y

% Construct main path
% Construct main path
        \pgfpathmoveto{\pgfpoint{\pgf@xa}{\pgf@ya+5pt}}
        \pgfpathlineto{\pgfpoint{\pgf@xa}{\pgf@yb}}
        \pgfpathlineto{\pgfpoint{\pgf@xb}{\pgf@yb}}
        \pgfpathlineto{\pgfpoint{\pgf@xb}{\pgf@ya+5pt}}
        \pgfpatharc{0}{-90}{5.0pt}
        \pgfpathlineto{\pgfpoint{\pgf@xa+5pt}{\pgf@ya}}
        \pgfpatharc{270}{180}{5.0pt}
%        \pgfpathcurveto{\pgfpoint{\pgf@xc}{\pgf@ya}\pgfpoint{\pgf@xa}{\pgf@ya}\pgfpoint{\pgf@xa}{\pgf@yc}}
        \pgfpathclose
    }
}

\pgfdeclareshape{howbecauseheadingbox}{%
    \inheritsavedanchors[from=rectangle]
    \inheritanchorborder[from=rectangle]

    \inheritanchor[from=rectangle]{center}
    \inheritanchor[from=rectangle]{south}
    \inheritanchor[from=rectangle]{west}
    \inheritanchor[from=rectangle]{north}
    \inheritanchor[from=rectangle]{east}
    \inheritanchor[from=rectangle]{base}
    \inheritanchor[from=rectangle]{north west}
    \inheritanchor[from=rectangle]{north east}
    \inheritanchor[from=rectangle]{south east}
    \inheritanchor[from=rectangle]{south west}

    \backgroundpath{%
        \southwest \pgf@xa=\pgf@x \pgf@ya=\pgf@y
        \northeast \pgf@xb=\pgf@x \pgf@yb=\pgf@y

% Construct main path
% Construct main path
        \pgfpathmoveto{\pgfpoint{\pgf@xa}{\pgf@ya}}
        \pgfpathlineto{\pgfpoint{\pgf@xa}{\pgf@yb}}
        \pgfpathlineto{\pgfpoint{\pgf@xb}{\pgf@yb}}
        \pgfpatharc{90}{0}{7.0pt}
        \pgfpathlineto{\pgfpoint{\pgf@xb+7pt}{\pgf@ya+7.0pt}}
        \pgfpatharc{0}{-90}{7.0pt}
        \pgfpathclose
    }
}


\newtcolorbox{openningbox}{
  enhanced,
%  jigsaw,
  opacityframe=.3,
  opacityfill=.3,
  width=80mm,
  left=3mm,
  right=3mm,
  top=3mm,
  bottom=.5mm,
  colback=white,
  colframe=white,
  arc=3mm,
  sharp corners=uphill,
    % frame empty,
    % arc = 8pt,
    lower separated = false,
    before upper = \tcbsubtitle{O QU\^E?},
    before lower = \tcbsubtitle{POR QU\^E?},
    fontupper = \sffamily,
    fontlower = \sffamily,
    sharp corners = north,
    subtitle style = {
      frame empty,
      fontupper = \fontsize{13pt}{0}\selectfont\titlefont\bfseries,
      colback = cor1,
      coltext = black!70,
      width = 3.5cm,
      arc = 8pt,
      rounded corners = east
    }
  }

\newcommand{\ChapNumTitle}[1]{%
  % \thispagestyle{plain}
  \thispagestyle{empty}
    \begin{tikzpicture}[remember picture, overlay]
        \coordinate (a) at ($(current page.north west) + (44.00mm,-44.00mm)$);
        \coordinate (c) at ($(current page.north west) + (148.54mm,-64.82mm)$);
        \coordinate (d) at ($(current page.north) + (0.00mm,-78.57mm)$);

        \begin{pgfonlayer}{top}
        \filldraw [fill=cor1, draw=cor1] ($(a) + (0.00mm,-25.00mm)$) -- ++(22.00mm,0.0mm) arc [start angle=270, end angle=360, x radius = 3.00mm, y
        radius=3.00mm] coordinate (c1) -- ++(0.00mm,44.00mm) coordinate (c2) arc [start angle=0, end angle=90, x radius=3.00mm, y radius=3.00mm] --
        ++(-44.00mm,0.00mm) arc [start angle=90, end angle=180, x radius=3.00mm, y radius=3.00mm] -- ++(0.00mm,-44.00mm) arc [start angle=180, end
        angle=270, x radius=3.00mm, y radius=3.00mm] -- cycle;
        \node[font=\chapnumfont, text=cor2] at (a) {\chapnamefont\thechapter};

        \coordinate (b) at ($(c1)!0.5!(c2)$);
        \node (title) at (b) [right, font=\chaptitlefont, text=cor2] {\begin{varwidth}{12cm}\baselineskip=1cm #1\end{varwidth}};
        % solution to the box O QUE and POR QUE?
        \node [right] at ($(current page.north)+(-15mm,-150mm)$) %
        {\begin{openningbox}
            \chapterwhat

            \tcblower

            \chapterbecause

            \end{openningbox}};


        \end{pgfonlayer}


        \begin{pgfonlayer}{bottom}
            \node[opacity=0.3,inner sep=0pt] at (current page.center){\includegraphics[width=\paperwidth,height=\paperheight]{\chapterillustration}};
        \end{pgfonlayer}

        \begin{pgfonlayer}{top}

        \end{pgfonlayer}
    \end{tikzpicture}

}

\newcommand{\ChapNoNumTitle}[1]{%
    \thispagestyle{plain}
    \begin{tikzpicture}[remember picture, overlay]
        \coordinate (a) at ($(current page.north west) + (44.00mm,-44.00mm)$);

        \filldraw [fill=cor1] ($(a) + (0.00mm,-25.00mm)$) -- ++(22.00mm,0.0mm) arc [start angle=270, end angle=360, x radius = 3.00mm, y
        radius=3.00mm] coordinate (c1) -- ++(0.00mm,44.00mm) coordinate (c2) arc [start angle=0, end angle=90, x radius=3.00mm, y radius=3.00mm] --
        ++(-44.00mm,0.00mm) arc [start angle=90, end angle=180, x radius=3.00mm, y radius=3.00mm] -- ++(0.00mm,-44.00mm) arc [start angle=180, end
        angle=270, x radius=3.00mm, y radius=3.00mm] -- cycle;
%        \node[textbf, text=cor2] at (a) {\chapnamefont\thechapter};

        \coordinate (b) at ($(c1)!0.5!(c2)$);
        \node (title) at (b) [right, font=\chaptitlefont] {#1};
    \end{tikzpicture}
}
\newif\ifnumberedchap
\numberedchaptrue

% Page number command
\newcommand{\leftpagenumber}{%
    \begin{tikzpicture}[remember picture, overlay]
        \coordinate (bottom) at ($(current page.south west) + (12mm,6.80mm)$);
            \begin{pgfonlayer}{main}
                % Square for page number
                \filldraw [fill=cor1, draw=cor1] (bottom) rectangle ++(-8.00mm,8.00mm) coordinate (corner);
            \end{pgfonlayer}

            \begin{pgfonlayer}{foreground}
                \node [text=white, font=\bfseries](page_number) at ($(bottom)!0.5!(corner)$) {\thepage};
            \end{pgfonlayer}
    \end{tikzpicture}
}

\newcommand{\rightpagenumber}{%
    \begin{tikzpicture}[remember picture, overlay]
        \coordinate (bottom) at ($(current page.south east) + (-12mm,6.80mm)$);
            \begin{pgfonlayer}{main}
                % Square for page number
                \filldraw [fill=cor1, draw=cor1] (bottom) rectangle ++(8.00mm,8.00mm) coordinate (corner);
            \end{pgfonlayer}

            \begin{pgfonlayer}{foreground}
                \node [text=white, font=\bfseries](page_number) at ($(bottom)!0.5!(corner)$) {\thepage};
            \end{pgfonlayer}

    \end{tikzpicture}
}

\newcommand{\drawpage}{%
    \begin{tikzpicture}[remember picture, overlay]
        \draw [green, thick] (current page.north west) -- (current page.north east) -- (current page.south east) -- (current page.south west) -- cycle;
    \end{tikzpicture}
}


% Draw headers. This places the tabs on the side of the page
\newcommand{\leftpagetab}{%
    \begin{tikzpicture}[remember picture, overlay]
        % \filldraw [fill=cor1, draw=cor1] (current page.north east) rectangle ($(current page.south east) + (-18.00mm,0.00mm)$) coordinate (bottom) {};
        \coordinate (bottom) at ($(current page.south west) + (18.00mm,0.00mm)$);

        % Draw color block
        \filldraw [fill=cor2, draw=cor2] ($(current page.north east) + (-12.00mm,0.00mm)$) rectangle ++(8.00mm,-8.00mm) coordinate (taba) {};
        % \filldraw [fill=cor1, draw=cor1] ($(current page.north east) + (-18.00mm,0.00mm)$) -- ++(0.00,-48.00mm) arc [start angle=180, end angle=90, x radius=2.00mm, y radius=2.00mm] -- ++(5.00mm,0.00mm) arc [start angle=270, end angle=360, x radius=2.00mm, y radius=2.00mm] |- cycle;

        % Square for chapter number
        \coordinate (tabb) at ($(taba) + (-8.00mm,-8.00mm)$);
        \coordinate (tabc) at ($(tabb) + (4.00mm,0.00mm)$);
        \node [fill=cor1, text=white, minimum height=8.00mm, minimum width=8.00mm] (chapternumber) at ($(tabb)!0.5!(taba)$) {\Large\thirdtitlefont\thechapter};
        % \filldraw [fill=cor2, draw=cor2] ($(current page.north east) + (-18.00mm,-26.00mm)$) -- ++(9.00mm,0.00mm) -- ++(0.00mm,-9.00mm) -| cycle;

		\coordinate (chaptertitle) at ($(taba) + (0.00mm,-25.00mm)$);

		
		\coordinate (chaptertitle2) at ($(chaptertitle) + (-4.00mm,0.00mm)$);
        % Add end tab
        \filldraw [fill=cor2, draw=cor2] ($(chaptertitle) + (-8.00mm,-4.00mm)$) rectangle ++(2.00mm,4.00mm);
        \filldraw [fill=white, draw=white] ($(chaptertitle) + (-6.00mm,-3.8mm)$) circle (5pt);
        \filldraw [fill=cor2, draw=cor2] ($(chaptertitle) + (-6.00mm,-1.7mm)$) rectangle ++(4.00mm,1.70mm);
        \filldraw [fill=cor2, draw=cor2] ($(chaptertitle) + (-2.00mm,-0.00mm)$) circle (5pt);



        % Add chapter name
        \node [fill=cor2, text=white] at (chaptertitle) [rotate=-90, below left, minimum height = 8.00mm] {\tiny\titlefont \MakeUppercase{\chaptername}};
		
        
        %\fill [fill=cor2] (chaptertitle2) -- ++(-4.00mm,0.00mm) coordinate (t1) arc [start angle=180, end angle=270, x radius=2.00mm, y radius=2.00mm]  -- ++(4.00mm,0.00mm) arc [start angle=270, end angle=360, x radius=2.00mm, y radius=2.00mm] coordinate (t2) -| cycle;

        % Add chapter title 
        \node  at (chaptertitle2) [right, yshift=-2mm, rotate=-90, text=black] {\tiny \pagetitlefont \leftmark};
        
        % Add right line
        \coordinate (innertop) at ($(current page.north east) + (-148mm,-16.00mm)$);
        \coordinate (innerbottom) at ($(current page.south east) + (-148mm,16.00mm)$);
        \coordinate (outertop) at ($(current page.north east) + (-18mm,-16.00mm)$);
        \coordinate (outerbottom) at ($(current page.south east) + (-18mm,16.00mm)$);
        \draw (innertop) --  (innerbottom);
        \draw (innertop) --  (outertop);
        \draw (outertop) --  (outerbottom);
        \draw (outerbottom) --  (innerbottom);

        % Water Mark

		% \coordinate (WMark) at ($(current page.east) - (122.00mm,0.00mm)$);
		% \fontsize{110}{130}\selectfont
  %       \node[text=grafite!20,rotate=-90] at (WMark) {Professor(a)};
    \end{tikzpicture}
}




\newcommand{\rightpagetab}{%
    \begin{tikzpicture}[remember picture, overlay]
        % \filldraw [fill=cor1, draw=cor1] (current page.north west) rectangle ($(current page.south west) + (18.00mm,0.00mm)$) coordinate (bottom) {};
        \coordinate (bottom) at ($(current page.south west) + (18.00mm,0.00mm)$);

        % Draw color block
        \filldraw [fill=cor2, draw=cor2] ($(current page.north west) + (12.00mm,0.00mm)$) rectangle ++(-8.00mm,-8.00mm) coordinate (taba) {};

        % Square for chapter number
        \coordinate (tabb) at ($(taba) + (8.00mm,-8.00mm)$);
        %%% \coordinate (tabc) at ($(tabb) + (4.00mm,0.00mm)$);
        \node [fill=cor1, text=white, minimum height=8.00mm, minimum width=8.00mm] (chapternumber) at ($(tabb)!0.5!(taba)$) {\Large\thirdtitlefont\thechapter};

        % Add chapter name
        \node [fill=cor2, text=white] (chaptertitle) at (tabb) [rotate=90, above left, minimum height = 8.00mm] {\tiny\titlefont \MakeUppercase{\chaptername}};

        % Add end tab
        \filldraw [fill=cor2, draw=cor2] (chaptertitle.south west) -- ++(0.00mm,-4.00mm) coordinate (t1) arc [start angle=0, end angle=90, x radius=2.00mm, y radius=2.00mm] -- ++(-4.00mm,0.00mm) arc [start angle=270, end angle=180, x radius=2.00mm, y radius=2.00mm] coordinate (t2) -| cycle;

        % Add chapter title
        \node (chaptertitle) at ($(t1)!0.5!(t2)$) [left, rotate=90, text=black, xshift=-0.00mm] {\tiny\pagetitlefont\leftmark};

        % Square for page number
        % \filldraw [fill=cor1, draw=cor1] ($(bottom) + (0.00mm, 20.00mm)$) rectangle ++(8.00mm,8.00mm);
        
        % Add right line
        \coordinate (innertop) at ($(current page.north west) + (148mm,-16.00mm)$);
        \coordinate (innerbottom) at ($(current page.south west) + (148mm,16.00mm)$);
        \coordinate (outertop) at ($(current page.north west) + (18mm,-16.00mm)$);
        \coordinate (outerbottom) at ($(current page.south west) + (18mm,16.00mm)$);
        \draw (innertop) --  (innerbottom);
        \draw (innertop) --  (outertop);
        \draw (outertop) --  (outerbottom);
        \draw (outerbottom) --  (innerbottom);


		\coordinate (WMark) at ($(current page.west) + (122.00mm,0.00mm)$);


		% \fontsize{110}{130}\selectfont
  %       \node[text=grafite!20,rotate=90] at (WMark) {Professor(a)};      
        
    \end{tikzpicture}
}

\let\oldnewline\newline
% redefinition: \newline is \relax
\newcommand\AnnCL{%
  \def\newline{\relax}}
% recovering the original meaning of \newline
\newcommand\RecCL{%
  \def\newline{\oldnewline}}

% Draw Headings






%% Exploring Headins

%% Explorando

\titleclass{\exploresec}{straight}[\section]
\newcounter{exploresec}
% \counterwithin{exploresec}{chapter}
\renewcommand*{\theexploresec}{\arabic{exploresec}}


\titleformat{\exploresec}[hang]{\large\color{session1}\secondtitlefont}{\theexploresec}{0pt}{\phantomsection
\begin{tikzpicture}[remember picture, scale=1.3]
        \filldraw [fill=session1, draw=session1] (-3.00mm,0.00mm) coordinate (a) -- ++(30.00mm,0.00mm) arc [start angle=90, end angle=0, x radius=2.00mm, y radius=2.00mm] coordinate (a1) -- ++(0.00mm,-2.00mm) coordinate (b1) arc [start angle=180, end angle=270, x radius=2.00mm, y radius=2.00mm] coordinate (b) -- ++(-30.00mm, 0.00mm) arc [start angle=270, end angle=180, x radius=2.00mm, y radius=2.00mm] -- ++(0.00mm,2.00mm) arc [start angle=0, end angle=90, x radius=2.00mm, y radius=2.00mm] -- cycle;
        \node (explore) at ($(a)!0.5!(b)$) [white] {\titlefont EXPLORANDO};
        \draw [session1, thick] (b) --++ ($(\textwidth,0.00mm)+(-72 mm,0 mm)$); % changed
        \node (title) at ($(a1)!0.5!(b1)$) [right] {\MakeUppercase{#1}};
\end{tikzpicture}}

\titlespacing*{\exploresec}{0pt}{8mm}{3pt}


\newcommand{\explore}[1]
{	\pdfbookmark[1]{#1}{}
	\needspace{\textheight}\checkoddpage\ifoddpage\else\cleardoublepage\fi%
	\stepcounter{exploresec}
	\exploresec{\texorpdfstring{}{Explorando: }#1}
	\def\currentcolor{session1}
	\setlist[enumerate]{label=\protect\bfseries\textcolor{\currentcolor}{\alph*)}}
	% choose secondary (depth 2) enumerate style..
	\setlist[enumerate,2]{label=\protect\bfseries\textcolor{\currentcolor}{\roman*)},ref=\protect\bfseries\textcolor{\currentcolor}{\alph*)}} 
}



%Praticando

\titleclass{\practicesec}{straight}[\section]
\newcounter{practicesec}
\renewcommand*{\thepracticesec}{\arabic{practicesec}}

\titleformat{\practicesec}[hang]{\large\color{session2}\secondtitlefont}{\thepracticesec}{0pt}{\phantomsection
% \needspace{\textheight}\checkoddpage\ifoddpage\else\cleardoublepage\fi%
\begin{tikzpicture}[remember picture, scale=1.3]% overlay removed from original, otherwise the stays behind the text, scale.
        \filldraw [fill=session2, draw=session2] (-3.00mm,0.00mm) coordinate (a) -- ++(30.00mm,0.00mm) arc [start angle=90, end angle=0, x radius=2.00mm, y radius=2.00mm] coordinate (a1) -- ++(0.00mm,-2.00mm) coordinate (b1) arc [start angle=180, end angle=270, x radius=2.00mm, y radius=2.00mm] coordinate (b) -- ++(-30.00mm, 0.00mm) arc [start angle=270, end angle=180, x radius=2.00mm, y radius=2.00mm] -- ++(0.00mm,2.00mm) arc [start angle=0, end angle=90, x radius=2.00mm, y radius=2.00mm] -- cycle;
        \node (explore) at ($(a)!0.5!(b)$) [white] {\titlefont PRATICANDO};
        \draw [session2, thick] (b) --++ ($(\textwidth,0.00mm)+(-72 mm,0 mm)$); % changed
        \node (title) at ($(a1)!0.5!(b1)$) [right] {\MakeUppercase{#1}};
\end{tikzpicture}}

\titlespacing*{\practicesec}{0pt}{8mm}{3pt}


\newcommand{\practice}[1]
{
	\needspace{.15\textheight}
	\stepcounter{practicesec}
	\practicesec{\texorpdfstring{}{Praticando: }#1}
	% \pdfbookmark[\thepracticesec]{Praticando: #1}{}
	\def\currentcolor{session2}
	\setlist[enumerate]{label=\protect\bfseries\textcolor{\currentcolor}{\alph*)}}
	% choose secondary (depth 2) enumerate style..
	\setlist[enumerate,2]{label=\protect\bfseries\textcolor{\currentcolor}{\roman*)},ref=\protect\bfseries\textcolor{\currentcolor}{\alph*)}} 
}


%Para Saber+

\titleclass{\knowsec}{straight}[\chapter]
\newcounter{knowsec}
\renewcommand*{\theknowsec}{\arabic{knowsec}}

\titleformat{\knowsec}[hang]{\large\color{session3}\secondtitlefont}{\theknowsec}{0pt}{
\phantomsection
\begin{tikzpicture}[remember picture, scale=1.3]% overlay removed from original, otherwise the stays behind the text, scale.
        \filldraw [fill=session3, draw=session3] (-3.00mm,0.00mm) coordinate (a) -- ++(30.00mm,0.00mm) arc [start angle=90, end angle=0, x radius=2.00mm, y radius=2.00mm] coordinate (a1) -- ++(0.00mm,-2.00mm) coordinate (b1) arc [start angle=180, end angle=270, x radius=2.00mm, y radius=2.00mm] coordinate (b) -- ++(-30.00mm, 0.00mm) arc [start angle=270, end angle=180, x radius=2.00mm, y radius=2.00mm] -- ++(0.00mm,2.00mm) arc [start angle=0, end angle=90, x radius=2.00mm, y radius=2.00mm] -- cycle;
        \node (explore) at ($(a)!0.5!(b)$) [white] {\titlefont PARA SABER \Huge{+}};
        \draw [session3, thick] (b) --++ ($(\textwidth,0.00mm)+(-72 mm,0 mm)$); % changed
        \node (title) at ($(a1)!0.5!(b1)$) [right] {\MakeUppercase{#1}};
\end{tikzpicture}}

\titlespacing*{\knowsec}{0pt}{8mm}{3pt}


\newcommand{\know}[1]
{
	\needspace{.15\textheight}
	\stepcounter{knowsec}
	\knowsec{\texorpdfstring{}{Para Saber +: }#1}
	% \pdfbookmark[\theknowsec]{Para Saber +: #1}{}
	\def\currentcolor{session3}
	\setlist[enumerate]{label=\protect\bfseries\textcolor{\currentcolor}{\alph*)}}
	% choose secondary (depth 2) enumerate style..
	\setlist[enumerate,2]{label=\protect\bfseries\textcolor{\currentcolor}{\roman*)},ref=\protect\bfseries\textcolor{\currentcolor}{\alph*)}} 
}


%Organizando


\titleclass{\arrangesec}{straight}[\section]
\newcounter{arrangesec}
\renewcommand*{\thearrangesec}{\arabic{arrangesec}}

\titleformat{\arrangesec}[hang]{\large\color{session4}\secondtitlefont}{\thearrangesec}{0pt}{\phantomsection
% \needspace{\textheight}\checkoddpage\ifoddpage\else\cleardoublepage\fi%
\begin{tikzpicture}[remember picture, scale=1.3]% overlay removed from original, otherwise the stays behind the text, scale.
        \filldraw [fill=session4, draw=session4] (-3.00mm,0.00mm) coordinate (a) -- ++(30.00mm,0.00mm) arc [start angle=90, end angle=0, x radius=2.00mm, y radius=2.00mm] coordinate (a1) -- ++(0.00mm,-2.00mm) coordinate (b1) arc [start angle=180, end angle=270, x radius=2.00mm, y radius=2.00mm] coordinate (b) -- ++(-30.00mm, 0.00mm) arc [start angle=270, end angle=180, x radius=2.00mm, y radius=2.00mm] -- ++(0.00mm,2.00mm) arc [start angle=0, end angle=90, x radius=2.00mm, y radius=2.00mm] -- cycle;
        \node (explore) at ($(a)!0.5!(b)$) [white] {\titlefont ORGANIZANDO};
        \draw [session3, thick] (b) --++ ($(\textwidth,0.00mm)+(-72 mm,0 mm)$); % changed
        \node (title) at ($(a1)!0.5!(b1)$) [right] {\MakeUppercase{#1}};
\end{tikzpicture}}

\titlespacing*{\arrangesec}{0pt}{8mm}{3pt}



\newcommand{\arrange}[1]
{
	\needspace{.15\textheight}
	\stepcounter{arrangesec}
	\arrangesec{\texorpdfstring{}{Organizando: }#1}
	% \pdfbookmark[\thearrangesec]{Organizando: #1}{}
	\def\currentcolor{session4}
	\setlist[enumerate]{label=\protect\bfseries\textcolor{\currentcolor}{\alph*)}}
	% choose secondary (depth 2) enumerate style..
	\setlist[enumerate,2]{label=\protect\bfseries\textcolor{\currentcolor}{\roman*)},ref=\protect\bfseries\textcolor{\currentcolor}{\alph*)}} 
}


\titleclass{\exercisesec}{straight}[\chapter]
\newcounter{exercisesec}
\renewcommand*{\theexercisesec}{\arabic{exercisesec}}


\titleformat{\exercisesec}[hang]{\large\color{cor1}\secondtitlefont}{\theexercisesec}{0pt}{\phantomsection
% \needspace{\textheight}\checkoddpage\ifoddpage\else\cleardoublepage\fi%
\begin{tikzpicture}[remember picture]
	\draw [draw=cor2, very thick] (-3.00mm,0.00mm) -- ++(0.00mm,2.00mm) -- ++(\textwidth , 0.00mm) -- ++(0mm,-2mm);% line
	\fill [draw=cor2, very thick , fill=cor1] ($(\textwidth/2 , 0mm) + (19mm,6mm)$) -- ++(0.00mm,-5.5 mm) coordinate (line1) arc [start angle=0, end angle=-90, x radius=2.00mm, y radius=2.00mm] -| ++(-38.00mm , 0.00mm)  -- ++ (0mm, 5.5 mm) coordinate (line2) arc [start angle=180, end angle=90, x radius=2.00mm, y radius=2.00mm] -- cycle;
	\node (exercise) at ($(line1)!0.5!(line2) + (0 mm, .8 mm)$) [text=white, font=\thirdtitlefont\Large] {EXERCÍCIOS};%\fontsize{13pt}{0}
\end{tikzpicture}}

\titlespacing*{\exercisesec}{0pt}{8mm}{1mm}



\newcommand{\exercise}[1]
{	\needspace{.15\textheight}
	\stepcounter{exercisesec}
	\exercisesec{\texorpdfstring{}{Exercícios}}
	% \pdfbookmark[\theexercisesec]{Exercícios}{}
	\def\currentcolor{cor1}
	\setlist[enumerate]{label=\protect\enumbox{\arabic*}}
	\setlist[enumerate,2]{label=\textit{\alph*)}
	% ,ref=\theenumi.\emph{\arabic*}
	}
	\setlist[enumerate,3]{label=\textit{\roman*)}
	% ,ref=\theenumi.\emph{\roman*}
	}
}

% \newcommand{\explore}[1]{\needspace{\textheight}
% \checkoddpage\ifoddpage\else\cleardoublepage\fi%
% \vspace{8 mm}
% \phantomsection\pdfbookmark[explore]{Explorando: {#1}}{#1}
% 	\noindent \begin{tikzpicture}[remember picture, scale=1.3]% overlay removed from original, otherwise the stays behind the text, scale.
%         \filldraw [fill=session1, draw=session1] (-3.00mm,0.00mm) coordinate (a) -- ++(30.00mm,0.00mm) arc [start angle=90, end angle=0, x radius=2.00mm, y radius=2.00mm] coordinate (a1) -- ++(0.00mm,-2.00mm) coordinate (b1) arc [start angle=180, end angle=270, x radius=2.00mm, y radius=2.00mm] coordinate (b) -- ++(-30.00mm, 0.00mm) arc [start angle=270, end angle=180, x radius=2.00mm, y radius=2.00mm] -- ++(0.00mm,2.00mm) arc [start angle=0, end angle=90, x radius=2.00mm, y radius=2.00mm] -- cycle;
%         \node (explore) at ($(a)!0.5!(b)$) [white] {\titlefont  EXPLORANDO};
%         \draw [session1, thick] (b) --++ ($(\textwidth,0.00mm)+(-72 mm,0 mm)$); % changed
%         \node (title) at ($(a1)!0.5!(b1)$) [right] {\large\secondtitlefont\textcolor{session1}{\MakeUppercase{ #1}}};
%     \end{tikzpicture}
%     \vspace{3pt}
    
%     	\def\currentcolor{session1}
% 	%setting depth 1 enumerate style
%         \setlist[enumerate]{label=\protect\bfseries\textcolor{\currentcolor}{\alph*)}}
% 	% choose secondary (depth 2) enumerate style..
% 	\setlist[enumerate,2]{label=\protect\bfseries\textcolor{\currentcolor}{\roman*)},ref=\protect\bfseries\textcolor{\currentcolor}{\alph*)}} 

    
% }

% \newcommand{\practice}[1]{%
% \vspace{8 mm}
% \needspace{.15\textheight}
% \phantomsection\pdfbookmark[practice]{Praticando: {#1}}{#1}

% \noindent     \begin{tikzpicture}[remember picture, scale=1.3]
%         \filldraw [fill=session2, draw=session2] (-3.00mm,0.00mm) coordinate (a) -- ++(30.00mm,0.00mm) arc [start angle=90, end angle=0, x radius=2.00mm, y radius=2.00mm] coordinate (a1) -- ++(0.00mm,-2.00mm) coordinate (b1) arc [start angle=180, end angle=270, x radius=2.00mm, y radius=2.00mm] coordinate (b) -- ++(-30.00mm, 0.00mm) arc [start angle=270, end angle=180, x radius=2.00mm, y radius=2.00mm] -- ++(0.00mm,2.00mm) arc [start angle=0, end angle=90, x radius=2.00mm, y radius=2.00mm] -- cycle;
%         \node (practice) at ($(a)!0.5!(b)$) [white] {\titlefont PRATICANDO};
%         \draw [session2, thick] (b) -- ++($(\textwidth,0.00mm)+(-72 mm,0 mm)$);
%         \node (title) at ($(a1)!0.5!(b1)$) [right] { \large\secondtitlefont\textcolor{session2}{\MakeUppercase{ #1}}};
%     \end{tikzpicture}
% \vspace{3pt}

% 	\def\currentcolor{session2}
% 	%setting depth 1 enumerate style
%         \setlist[enumerate]{label=\protect\bfseries\textcolor{\currentcolor}{\alph*)}}
% 	% choose secondary (depth 2) enumerate style..
% 	\setlist[enumerate,2]{label=\protect\bfseries\textcolor{\currentcolor}{\roman*)},ref=\protect\bfseries\textcolor{\currentcolor}{\alph*)}} 


% }

% \newcommand{\know}[1]{%
% \vspace{8 mm}
% \needspace{.15\textheight}
% \phantomsection\pdfbookmark[know]{Para Saber +: {#1}}{#1}

% \noindent     \begin{tikzpicture}[remember picture, scale=1.3]
%         \filldraw [fill=session3, draw=session3] (-3.00mm,0.00mm) coordinate (a) -- ++(30.00mm,0.00mm) arc [start angle=90, end angle=0, x radius=2.00mm, y radius=2.00mm] coordinate (a1) -- ++(0.00mm,-2.00mm) coordinate (b1) arc [start angle=180, end angle=270, x radius=2.00mm, y radius=2.00mm] coordinate (b) -- ++(-30.00mm, 0.00mm) arc [start angle=270, end angle=180, x radius=2.00mm, y radius=2.00mm] -- ++(0.00mm,2.00mm) arc [start angle=0, end angle=90, x radius=2.00mm, y radius=2.00mm] -- cycle;
%         \node (know) at ($(a)!0.5!(b)$) [white] {\titlefont PARA SABER \Huge{+}};
%         \draw [session3, thick] (b) -- ++($(\textwidth,0.00mm)+(-72 mm,0 mm)$);
%         \node (title) at ($(a1)!0.5!(b1)$) [right] {\large\secondtitlefont\textcolor{session3}{\MakeUppercase{ #1}}};
%     \end{tikzpicture}
% \vspace{3pt}

% 	\def\currentcolor{session3}
% 	%setting depth 1 enumerate style
%        \setlist[enumerate]{label=\protect\bfseries\textcolor{\currentcolor}{\alph*)}}
% 	% choose secondary (depth 2) enumerate style..
% 	\setlist[enumerate,2]{label=\protect\bfseries\textcolor{\currentcolor}{\roman*)},ref=\protect\bfseries\textcolor{\currentcolor}{\alph*)}} 


% }

% \newcommand{\arrange}[1]{%
% \vspace{8 mm}
% \needspace{.15\textheight}
% \phantomsection\pdfbookmark[arrange]{Organizando: {#1}}{#1}

% \noindent     \begin{tikzpicture}[remember picture, scale=1.3]
%         \filldraw [fill=session4, draw=session4] (-3.00mm,0.00mm) coordinate (a) -- ++(30.00mm,0.00mm) arc [start angle=90, end angle=0, x radius=2.00mm, y radius=2.00mm] coordinate (a1) -- ++(0.00mm,-2.00mm) coordinate (b1) arc [start angle=180, end angle=270, x radius=2.00mm, y radius=2.00mm] coordinate (b) -- ++(-30.00mm, 0.00mm) arc [start angle=270, end angle=180, x radius=2.00mm, y radius=2.00mm] -- ++(0.00mm,2.00mm) arc [start angle=0, end angle=90, x radius=2.00mm, y radius=2.00mm] -- cycle;
%         \node (arrange) at ($(a)!0.5!(b)$) [white] {\titlefont ORGANIZANDO};
%         \draw [session4, thick] (b) -- ++($(\textwidth,0.00mm)+(-72 mm,0 mm)$);
%         \node (title) at ($(a1)!0.5!(b1)$) [right] {\large\secondtitlefont\textcolor{session4}{\MakeUppercase{ #1}}\newline};
%     \end{tikzpicture}
% \vspace{3pt}

% 	\def\currentcolor{session4}
% 	%setting depth 1 enumerate style
%         \setlist[enumerate]{label=\protect\bfseries\textcolor{\currentcolor}{\alph*)}}
% 	% choose secondary (depth 2) enumerate style..
% 	\setlist[enumerate,2]{label=\protect\bfseries\textcolor{\currentcolor}{\roman*)},ref=\protect\bfseries\textcolor{\currentcolor}{\alph*)}} 


% }

% %% Exercise section
% \newcommand{\exercise}{
% \needspace{.15\textheight}
% \vspace{8 mm}

% 	\noindent \begin{tikzpicture}[remember picture]
% 	\draw [draw=cor2, very thick] (-3.00mm,0.00mm) -- ++(0.00mm,2.00mm) -- ++(\textwidth , 0.00mm) -- ++(0mm,-2mm);% line
% 	\fill [draw=cor2, very thick , fill=cor1] ($(\textwidth/2 , 0mm) + (19mm,6mm)$) -- ++(0.00mm,-5.5 mm) coordinate (line1) arc [start angle=0, end angle=-90, x radius=2.00mm, y radius=2.00mm] -| ++(-38.00mm , 0.00mm)  -- ++ (0mm, 5.5 mm) coordinate (line2) arc [start angle=180, end angle=90, x radius=2.00mm, y radius=2.00mm] -- cycle;
% 	\node (exercise) at ($(line1)!0.5!(line2) + (0 mm, .8 mm)$) [text=white, font=\thirdtitlefont\Large] {EXERCÍCIOS};%\fontsize{13pt}{0}
% 	\end{tikzpicture}
% 	\vspace{1mm}
%     \def\currentcolor{cor1}


% \setlist[enumerate]{label=\protect\enumbox{\arabic*}}
% \setlist[enumerate,2]{label=\textit{\alph*)}
% % ,ref=\theenumi.\emph{\arabic*}
% }
% \setlist[enumerate,3]{label=\textit{\roman*)}
% % ,ref=\theenumi.\emph{\roman*}
% }
% \phantomsection\pdfbookmark[exercise]{Exercícios}{}
% }


% Defining the section style for the diferent names that eventually appear.
\setsecnumdepth{chapter} % just number the chapter, not lower levels.

% Clear page after chapter title
\renewcommand\memendofchapterhook{%
    \clearpage\m@mindentafterchapter\@afterheading}

% Setup page
\makechapterstyle{livroaberto}{%
    \renewcommand{\chapnamefont}{}
    \renewcommand{\chapnumfont}{\titlefont\bfseries\fontsize{140pt}{0}\selectfont}
    \renewcommand{\chaptitlefont}{\titlefont\bfseries\fontsize{25pt}{0}\selectfont}
    \renewcommand{\chapternamenum}{}
    \renewcommand{\afterchapternum}{}
    \renewcommand{\printchaptername}{}
    \renewcommand{\printchapternum}{}
    \renewcommand{\printchapternonum}{\global\numberedchapfalse}
    \renewcommand{\printchaptertitle}[1]{%
    \thispagestyle{plain}
        \ifnumberedchap
            \ChapNumTitle{##1}
        \else
            \ChapNoNumTitle{##1}
        \fi
        \global\numberedchaptrue
        }
    \renewcommand{\afterchaptertitle}{\clearpage}
%    \setlength{\beforechapskip}{}
%    \setlength{\midchapskip}{}
%    \setlength{\afterchapskip}{}

}
\renewcommand{\chaptermarksn}[1]{}


\makepagestyle{plain}
\makeevenfoot{plain}{}{\leftpagenumber}{}
\makeoddfoot{plain}{}{\rightpagenumber}{}

\makepagestyle{livroaberto}

\makeevenfoot{livroaberto}{}{\leftpagenumber}{}
\makeoddfoot{livroaberto}{}{\rightpagenumber}{}
\makeevenhead{livroaberto}{}{\rightpagetab}{}
\makeoddhead{livroaberto}{}{\leftpagetab}{}


\chapterstyle{livroaberto}
\pagestyle{livroaberto}

%Seção 5 - Livro Aberto Lists

%***************************************************************************************
% This package creates the heading designs for the Livro Aberto textbook.
%
% This package was written and created by
% David Latchman, TeXnical Designs
% Contact: david.latchman@texnical-designs.com for problems
%
%

% \def\fileversion{v1.0}
% \def\filename{livroabertolists}
% \def\filetitle{Livro Aberto Lists package}
% \def\filedate{2018/02/13}
% \def\docdate {2018/02/13}


\NeedsTeXFormat{LaTeX2e}
\ProvidesPackage{\filename}

\typeout{Package: `\filetitle\space\filename\space\fileversion \@spaces <\filedate>'}
\typeout{English documentation as of <\docdate>}

%% ---------------------------------------------------------------------
%%
%% Loads the ifthen package
%%
\RequirePackage{ifthen}

\RequirePackage{tikz}
\usetikzlibrary{calc}
\usetikzlibrary{shapes, decorations}
\usetikzlibrary{intersections,through,backgrounds,fit}
\RequirePackage{tikzpagenodes}
\RequirePackage[inline]{enumitem}

%Alinha equações com item de listas
\newcommand\Item[1][]{%
  \ifx\relax#1\relax  \item \else \item[#1] \fi
  \abovedisplayskip=0pt\abovedisplayshortskip=0pt~\vspace*{-\baselineskip}}


\pgfdeclareshape{enumbox}{%
    \inheritsavedanchors[from=rectangle]
    \inheritanchorborder[from=rectangle]

    \inheritanchor[from=rectangle]{center}
    \inheritanchor[from=rectangle]{south}
    \inheritanchor[from=rectangle]{west}
    \inheritanchor[from=rectangle]{north}
    \inheritanchor[from=rectangle]{east}
    \inheritanchor[from=rectangle]{base}

    \backgroundpath{%
        \southwest \pgf@xa=\pgf@x \pgf@ya=\pgf@y
        \northeast \pgf@xb=\pgf@x \pgf@yb=\pgf@y

% Construct main path
        \pgfpathmoveto{\pgfpoint{\pgf@xa}{\pgf@ya+2pt}}
        \pgfpathlineto{\pgfpoint{\pgf@xa}{\pgf@yb}}
        \pgfpathlineto{\pgfpoint{\pgf@xb-2pt}{\pgf@yb}}
        \pgfpatharc{90}{0}{2.0pt}
        \pgfpathlineto{\pgfpoint{\pgf@xb}{\pgf@ya}}
        \pgfpathlineto{\pgfpoint{\pgf@xa+2pt}{\pgf@ya}}
        \pgfpatharc{270}{180}{2.0pt}
%        \pgfpathcurveto{\pgfpoint{\pgf@xc}{\pgf@ya}\pgfpoint{\pgf@xa}{\pgf@ya}\pgfpoint{\pgf@xa}{\pgf@yc}}
        \pgfpathclose
    }
}

\pgfdeclareshape{itembox}{%
    \inheritsavedanchors[from=rectangle]
    \inheritanchorborder[from=rectangle]

    \inheritanchor[from=rectangle]{center}
    \inheritanchor[from=rectangle]{south}
    \inheritanchor[from=rectangle]{west}
    \inheritanchor[from=rectangle]{north}
    \inheritanchor[from=rectangle]{east}
    \inheritanchor[from=rectangle]{base}

    \backgroundpath{%
        \southwest \pgf@xa=\pgf@x \pgf@ya=\pgf@y
        \northeast \pgf@xb=\pgf@x \pgf@yb=\pgf@y

% Construct main path
        \pgfpathmoveto{\pgfpoint{\pgf@xa}{\pgf@ya}}
        \pgfpathlineto{\pgfpoint{\pgf@xa}{\pgf@yb}}
        \pgfpathlineto{\pgfpoint{\pgf@xb}{\pgf@yb}}
        \pgfpathlineto{\pgfpoint{\pgf@xb}{\pgf@ya}}
        \pgfpathclose
    }
}

\newcommand*{\enumbox}[1]{%
    \tikz [baseline=(numbereditem.base)]{%
        \node[shape=enumbox,text=white, font=\bfseries\fontsize{11pt}{0}\selectfont, draw=black, 
        fill=cor1, 
        minimum width=6.00mm, minimum height=6.00mm] (numbereditem) {#1};
    }
}

\newcommand*{\@endnotebox}[1]{%
    \tikz [baseline=(numbereditem.base)]{%
        \node[shape=enumbox,text=linkscolor, font=\fontsize{11pt}{0}\selectfont\bfseries, 
        fill=box2!70, drop shadow,
        minimum width=4.00mm, minimum height=3.00mm,sharp corners] (numbereditem) {#1};
    }
}

\newcommand*{\notebox}[1]{%
    \tikz [baseline=(numbereditem.base)]{%
        \node[shape=enumbox,text=linkscolor, font=\fontsize{11pt}{0}\selectfont\bfseries, 
        fill=box2!70, drop shadow,
        minimum width=16.0mm, minimum height=3.00mm,sharp corners, align=left] (numbereditem) {#1};
    }
}

\newcommand*{\enumdefault}[1]{%
    \tikz [baseline=(numbereditem.base)]{%
        \node[shape=enumbox, font=\bfseries\fontsize{11pt}{0}\selectfont, draw=black, minimum width=6.00mm, minimum height=6.00mm] (numbereditem) {#1};
    }
}




% Defining an alternative shape for itemize, where the bullets are actually squares.
\newcommand*{\itembox}{%
    \tikz [baseline=(boxitem.base)] {%
        \node[shape=rectangle, draw=black, minimum width=6.00mm, minimum height=6.00mm] (boxitem) {};
    }
}

% Defining an alternative shape for itemize, where the bullets are actually colored squares.
\newcommand*{\itemcoloredsquare}{%
    \tikz {%
        \node[shape=rectangle, fill=\currentcolor, minimum width=1.50mm, minimum height=1.50mm] (boxitem) {};
    }
}


% Set up lists
\setlist[enumerate]{label=\protect\enumbox{\alph*)}}
%\setlist[itemize]{label=\protect\itembox}
\setlist[itemize]{label=\protect\itemcoloredsquare}


%Seção 6 - Livro Aberto Page

%***************************************************************************************
% This package sets up the page dimensions used in the Livro Aberto textbook.
%
% This package was written and created by
% David Latchman, TeXnical Designs
% Contact: david.latchman@texnical-designs.com for problems
%
%

% \def\fileversion{v1.0}
% \def\filename{livroabertopage}
% \def\filetitle{Livro Aberto page package}
% \def\filedate{2018/02/13}
% \def\docdate {2018/02/13}


\NeedsTeXFormat{LaTeX2e}
\ProvidesPackage{\filename}

\typeout{Package: `\filetitle\space\filename\space\fileversion \@spaces <\filedate>'}
\typeout{English documentation as of <\docdate>}

%% ---------------------------------------------------------------------
%%
%% Loads the ifthen package
%%


%% Set Page Size
% \setstocksize{11.00in}{8.50in}
\setstocksize{275mm}{204mm}
\settrimmedsize{275mm}{204mm}{*}
% \setstocksize{275mm}{204mm}
% \settrimmedsize{\stockheight}{\stockwidth}{*}


% \settrims{2.20mm}{5.95mm}

\settypeblocksize{241mm}{162mm}{*}
\setlrmarginsandblock{20mm}{22mm}{*}
\setulmarginsandblock{17mm}{17.00mm}{*}

% \setlrmargins{*}{1.25in}{*}
% \setulmargins{*}{*}{0.85}
% \setlrmarginsandblock{1.00in}{1.00in}{*}
\setcolsepandrule{8.00mm}{0mm}


% Set Paragraphs to no indent with a 6pt space between paragraphs
\setlength{\parindent}{0em}
 \setlength{\parskip}{6pt}

%% Modify Header and Footer spacing
% \setheadfoot{\onelineskip}{2\onelineskip}

\checkandfixthelayout

% Add some of the required graphics pacjakes to make things work
\RequirePackage{tikz}
\usetikzlibrary{calc}
\usetikzlibrary{shapes, decorations}
\usetikzlibrary{intersections,through,backgrounds}
\usetikzlibrary{positioning,fit}
\RequirePackage{tikzpagenodes}

\RequirePackage{environ}

% Bibliotecas e pacotes TikZ e Asymptote
\RequirePackage{tikz}
\usetikzlibrary{
    calc,trees,shadows,shadows.blur,positioning,arrows,chains,shapes.geometric, decorations, decorations.pathreplacing,decorations.pathmorphing,shapes,
    matrix,fit,tikzmark,intersections,through,backgrounds,arrows.meta,datavisualization}

\tikzstyle{vetor}=[->,line width=0.8pt,color=primario,>=latex]
\tikzstyle{ponto}=[circle, minimum size=2pt, inner sep=0, draw=black, fill=black, shift only, label={}]
\tikzstyle{every node}=[shift only]
\tikzstyle{eixos}=[->]
\tikzstyle{curvafun}=[very thick, color=primario]
\tikzstyle{grafica}+=[line cap=round,line join=round,color=terciario, very thick]
\newcommand{\vetoresmapaantartida}{
      \filldraw [white, fill opacity=100] (-4,-3) rectangle (4,3);
      \node at (-4,-3) {\includegraphics[width=1cm]{/home/cheque/Dropbox/Tikz-LivroAberto/Vetores/rosa.png}};
      \node [above, scale=0.5] at (-4,-2.5) {N};
      \node [left, scale=0.5] at (-4.5,-3) {E};
      \node [right, scale=0.5] at (-3.5,-3) {L};
      \node [below, scale=0.5] at (-4,-3.5) {S};
      \node [ponto] (A) at (0,0) {};
      \node [below] at (A) {$A$};
      \draw [|-|] (2,-3.5) -- (4,-3.5);
      \node [above, scale=0.5] at (3,-3.5) {10 km};
}

\RequirePackage{tikzpagenodes}

\RequirePackage{environ}

% Define asymptote picture preamble

\begin{asydef}
import three;
import graph3;
import solids;
import math;
import labelpath3;
import labelpath;
settings.render=16;
//attach=true;

pen azul = rgb(0,94/255,176/255); //azul
pen verde = rgb(75/255,173/255,49/255); //verde
pen vinho = rgb(183/255,13/255,40/255); //vinho
pen laranja = rgb(236/255,107/255,16/255); //laranja
pen agua = rgb(0,165/255,157/255); //verde agua
pen box2 = rgb(240/255,238/255,201/255);


texpreamble("\usepackage{fontspec}\setmainfont{Exo2-Light}");
\end{asydef}

% Define asymptote picture directory

\def\asydir{asytmp-perspectiva2}

% Define the layers to be used in document.
\pgfdeclarelayer{background}
\pgfdeclarelayer{foreground}
\pgfdeclarelayer{top}
\pgfdeclarelayer{bottom}
\pgfsetlayers{bottom,background,main,foreground,top}


%Seção 7 - Livro Aberto Professor
%***************************************************************************************
%
%This section configures the aluno document to be able to compile the professor document
%

%%%%%%%%%%%%%%%%%%%%%%%%%%%%%%%%%
%%%%%%%%%%%%%%%% Rotate title

\RequirePackage{soul}


\newlength\srotate@axis
\newlength\srotate@width
\DeclareRobustCommand*{\srotate}{%
  \SOUL@setup
  \let\SOUL@everyhyphen\srotate@everyhyphen
  \let\SOUL@everyexhyphen\srotate@everyexhyphen
  \let\SOUL@everytoken\srotate@everytoken
  \settoheight\srotate@axis{$\vcenter{}$}%
  \settoheight\srotate@width{\raisebox{\depth}{XQ}}%
  \SOUL@
}
\newcommand*{\srotate@everytoken}{%
  \srotate@token{\the\SOUL@token}%
}
\newcommand*{\srotate@everyexhyphen}[1]{%
  \srotate@token{\raisebox{\srotate@axis}[.1ex][.1ex]{#1}}%
}
\newcommand*{\srotate@everyhyphen}{%
  \discretionary{%
    \srotate@everyexhyphen{-}%
  }{}{}%
}
\newcommand*{\srotate@token}[1]{%
  \kern.02ex
  \hbox to \srotate@width{%
    \hfil
    \raisebox{\srotate@axis}{\rotatebox[origin=B]{90}{#1}}%
    \hfil
  }%
  \kern.1ex
}


%%%%%%%%%%%%%%% ENDNOTES

\RequirePackage{endnotes}

\def\enotesize{\normalsize}
%%% hyperendnotes.sty
\makeatletter
\newif\ifenotelinks
\newcounter{Hendnote}
% Redefining portions of endnotes-package:
\let\savedhref\href
\let\savedurl\url
\def\endnotemark{%
\@ifnextchar[\@xendnotemark{%
\stepcounter{endnote}%
\protected@xdef\@theenmark{\textbf{Nota \theendnote}}%
\protected@xdef\@theenvalue{\number\c@endnote}%
\@endnotemark
}%
}%
\def\@xendnotemark[#1]{%
\begingroup\c@endnote#1\relax
\unrestored@protected@xdef\@theenmark{\theendnote}%
\unrestored@protected@xdef\@theenvalue{\number\c@endnote}%
\endgroup
\@endnotemark
}%
\def\endnotetext{%
\@ifnextchar[\@xendnotenext{%
\protected@xdef\@theenmark{\theendnote}%
\protected@xdef\@theenvalue{\number\c@endnote}%
\@endnotetext
}%
}%

\def\@xendnotenext[#1]{%
\begingroup
\c@endnote=#1\relax
\unrestored@protected@xdef\@theenmark{\theendnote}%
\unrestored@protected@xdef\@theenvalue{\number\c@endnote}%
\endgroup
\@endnotetext
}%
\def\endnote{%
\@ifnextchar[\@xendnote{%
\stepcounter{endnote}%
\protected@xdef\@theenmark{\theendnote}%
\protected@xdef\@theenvalue{\number\c@endnote}%
\@endnotemark\@endnotetext
}%
}%
\def\@xendnote[#1]{%
\begingroup
\c@endnote=#1\relax
\unrestored@protected@xdef\@theenmark{\theendnote}%
\unrestored@protected@xdef\@theenvalue{\number\c@endnote}%
\show\@theenvalue
\endgroup
\@endnotemark\@endnotetext
}%
\def\@endnotemark{%
\leavevmode
\ifhmode
\edef\@x@sf{\the\spacefactor}\nobreak
\fi
\ifenotelinks
\expandafter\@firstofone
\else
\expandafter\@gobble
\fi
{%
\Hy@raisedlink{%
\hyper@@anchor{Hendnotepage.\@theenvalue}{\empty}%
}%
}%
\hyper@linkstart{link}{Hendnote.\@theenvalue}%
\makeenmark
\hyper@linkend
\ifhmode
\spacefactor\@x@sf
\fi
\relax
}%
\long\def\@endnotetext#1{%
\if@enotesopen
\else
\@openenotes
\fi
\immediate\write\@enotes{%
\@doanenote{\@theenmark}{\@theenvalue}%
}%
\begingroup
\def\next{#1}%
\newlinechar='40
\immediate\write\@enotes{\meaning\next}%
\endgroup
\immediate\write\@enotes{%
\@endanenote
}%
}%
\def\theendnotes{%
\immediate\closeout\@enotes
\global\@enotesopenfalse
\begingroup
\makeatletter
\edef\@tempa{`\string>}%
\ifnum\catcode\@tempa=12
\let\@ResetGT\relax
\else
\edef\@ResetGT{\noexpand\catcode\@tempa=\the\catcode\@tempa}%
\@makeother\>%
\fi
\def\@doanenote##1##2##3>{%
\def\@theenmark{##1}%
\def\@theenvalue{##2}%
\par
\smallskip %<-small vertical gap between endnotes
\begingroup
\def\href{\expandafter\savedhref}%
\def\url{\expandafter\savedurl}%
\@ResetGT
\edef\@currentlabel{\csname p@endnote\endcsname\@theenmark}%
\enoteformat
}%
\def\@endanenote{%
\par\endgroup
}%
% Redefine, how numbers are formatted in the endnotes-section:
\renewcommand*\@makeenmark{%
\normalfont\protect\@endnotebox{\@theenmark}~%
}%
% header of endnotes-section
% \enoteheading
% font-size of endnotes
\enotesize
\input{\jobname.ent}%
\endgroup
}%
\def\enoteformat{%
\rightskip\z@
\leftskip1.8em
\parindent\z@
\leavevmode\llap{%
\setcounter{Hendnote}{\@theenvalue}%
\addtocounter{Hendnote}{-1}%
\refstepcounter{Hendnote}%
\ifenotelinks
\expandafter\@secondoftwo
\else
\expandafter\@firstoftwo
\fi
{\@firstofone}%
{\hyperlink{Hendnotepage.\@theenvalue}}%
{\makeenmark}%
}%
}%
% stop redefining portions of endnotes-package:
\makeatother
% Toggle switch in order to turn on/off back-links in the
% endnote-section:
\enotelinkstrue
%\enotelinksfalse


\renewcommand{\notesname}{Notas}

\makeatletter
%\def\@makeenmark{\hbox{\@textsuperscript{\normalfont\@theenmark}}}
\def\@makeenmark{\hbox{\normalfont\@theenmark}}

\titlespacing*{\section}
{0pt}{2mm}{1mm}
\titlespacing*{\subsection}
{0pt}{1mm}{1mm}

\setitemize{label=\itemcoloredsquare,topsep=2pt, itemsep=1pt}


%%%%%%%%%%%%%%%%  PÁGINA DE TEXTO

\newcounter{mypresentation}
\setcounter{mypresentation}{0}


\newenvironment{apresentacao}[1]
{
\setcounter{mypresentation}{\value{page}}
\frontmatter
\setcounter{mypresentation}{\value{page}}
\def\currentcolor{cor1}
\newgeometry{textwidth=840pt,outer=50pt, bottom=45pt, top=46pt}
\pagestyle{plain}
\setlist[enumerate]{wide, label=\titem{\alph*)},itemsep=0pt,topsep=2.5pt}
\setlist[itemize]{wide, label=\textcolor{\currentcolor}{\small$\blacksquare$}, topsep=5pt} 
\section{#1}
\begin{multicols}{2}}
{\clearpage\end{multicols}
\pagebreak\restoregeometry 
\mainmatter
\def\currentcolor{session1}
\cleardoublepage
\setcounter{page}{\value{mypresentation}}
}

\newcounter{mypagecount}
\setcounter{mypagecount}{0}

\newenvironment{paginatexto}[1]
{
\renewcommand\frontmatter{%
  \cleardoublepage
  \@mainmatterfalse
  %\pagenumbering{roman}% Don't reset
}
\renewcommand\mainmatter{%
  \cleardoublepage
  \@mainmattertrue
  %\pagenumbering{arabic}% Don't reset
}\clearpage
\setcounter{mypagecount}{\value{page}}
\cleardoublepage
\setcounter{mypagecount}{\value{page}}
\def\currentcolor{session1}
\newgeometry{textwidth=840pt,outer=50pt, bottom=45pt, top=46pt}
\thispagestyle{empty}\pagestyle{empty}
\setlist[enumerate]{wide, label=\titem{\alph*)},itemsep=0pt,topsep=2.5pt}
\setlist[itemize]{wide, label=\textcolor{\currentcolor}{\small$\blacksquare$}, topsep=5pt}
\section{#1}
\begin{multicols}{2}}
{\clearpage\end{multicols}\def\currentcolor{session1}\pagebreak\restoregeometry\cleardoublepage\setcounter{page}{\value{mypagecount}}
}



%%%%%%%%%%%%%%%%  PÁGINA DE NOTAS



\newenvironment{paginanotas}
{\newgeometry{textwidth=840pt,outer=50pt}\pagestyle{empty} \begin{multicols}{4}}
{\clearpage\end{multicols}\pagebreak\restoregeometry\mainmatter}


%%%%%%%%%%%%% CAIXA DE TEXTO

 
\newcommand\semihuge{\@setfontsize\semiHuge{16.72}{21.38}}


\newenvironment{texto}[1]
{\marginpar{
  \setlength\parskip{6pt}
	\vspace{.05em}
  \setlist[enumerate]{label=OE\arabic*,wide}
  \setlist[itemize,1]{wide, label=\textcolor{\currentcolor}{\small$\blacksquare$}, topsep=5pt}
 
	#1}
}



%%%%%%%%%%%%% CAIXA DE OBJETIVO ESPECÍFICO

\newenvironment{objectives}[4]
{\ifnum #3=1   
{\ifnum #4=1

\marginpar{\vspace{.05em}\begin{tcolorbox}[blanker, enhanced, enforce breakable, before skip = 5mm, after skip = 5mm, colframe=\currentcolor, left=3mm, top=6mm,segmentation style={draw=none,fill=none}, bottom=3mm, colback=white, %colframe=box1,
borderline east={2mm}{0pt}{white}, toprule=1pt, bottomrule=1pt, leftrule=1pt, rightrule=-1pt, outer arc = 3mm, arc = 3mm,
fonttitle=\bfseries\large,  % Should be 12 pt. Don't know how to set this.
title= Objetivos Específicos, 
boxed title style={empty,arc=0pt,outer arc=0pt,boxrule=0pt}, 
attach boxed title to top right={yshift=.5mm-.8pt},
underlay boxed title={
\draw[line width=3pt,white] ($(frame.north east)+(0,-1pt)$)--($(frame.north east)+(-15pt,-1pt)$);
\filldraw [\currentcolor] ($(title.north west) + (0 mm,-3.50mm)$) -- ++(0.00mm,2.5mm) arc [start angle=180, end angle=90, y radius=2.00mm, x radius=2.00mm] -| ($(title.north east)+(0,-5mm)$) arc [start angle=0, end angle=-90, x radius=2.00mm, y radius=2.00mm] -| cycle;
}, 
before upper={\setlength{\parskip}{6pt}\vspace{-15pt}\tcbsubtitle{#1}
\normalfont\setlist[enumerate]{label=OE\arabic*,wide}\setlist[itemize,1]{wide, label=\textcolor{\currentcolor}{\small$\blacksquare$}, topsep=5pt}},
% {\boxtitlefont\setlist[enumerate]{label=OE\arabic*,wide}\setlist[itemize,1]{wide, label=\textcolor{penaqua}{\small$\blacksquare$}}
% {\vspace{1em}\large\textbf{Levar o aluno a}}\normalsize},
    subtitle style = {
      frame empty,
      fontupper = \bfseries\large,
      colback = white,
      coltext = \currentcolor,
      arc = 8pt,
      top=0pt,
      bottom=0pt
    },
    middle=5pt,
    break at=-\baselineskip/0pt] #2 \end{tcolorbox}}
    

\else

\marginpar{\vspace{.05em}\begin{tcolorbox}[blanker, enhanced, enforce breakable, before skip = 5mm, after skip = 5mm, colframe=\currentcolor, right=3mm, top=6mm,segmentation style={draw=none,fill=none}, bottom=3mm, colback=white, %colframe=box1,
borderline west={2mm}{0pt}{white}, toprule=1pt, bottomrule=1pt, rightrule=1pt, leftrule=-1pt, outer arc = 3mm, arc = 3mm,
fonttitle=\bfseries\large,  % Should be 12 pt. Don't know how to set this.
title= Objetivos Específicos, 
boxed title style={empty,arc=0pt,outer arc=0pt,boxrule=0pt}, 
attach boxed title to top left={yshift=.5mm-.8pt},
underlay boxed title={
\draw[line width=3pt,white] ($(frame.north west)+(0,-1pt)$)--($(frame.north west)+(15pt,-1pt)$);
\filldraw [\currentcolor] ($(title.north west) + (0 mm,-3.50mm)$) -- ++(0.00mm,2.5mm) arc [start angle=180, end angle=90, y radius=2.00mm, x radius=2.00mm] -| ($(title.north east)+(0,-5mm)$) arc [start angle=0, end angle=-90, x radius=2.00mm, y radius=2.00mm] -| cycle;
}, 
before upper={\setlength{\parskip}{6pt}\vspace{-15pt}\tcbsubtitle{#1}
\normalfont\setlist[enumerate]{label=OE\arabic*,wide}\setlist[itemize,1]{wide, label=\textcolor{\currentcolor}{\small$\blacksquare$}, topsep=5pt}},
% {\boxtitlefont\setlist[enumerate]{label=OE\arabic*,wide}\setlist[itemize,1]{wide, label=\textcolor{penaqua}{\small$\blacksquare$}}
% {\vspace{1em}\large\textbf{Levar o aluno a}}\normalsize},
    subtitle style = {
      frame empty,
      fontupper = \bfseries\large,
      colback = white,
      coltext = \currentcolor,
      arc = 8pt,
      top=0pt,
      bottom=0pt
    },
    middle=5pt,
    break at=-\baselineskip/0pt] #2 \end{tcolorbox}}
    
\fi}


\else

{\ifnum #4=9

\marginpar{\notebox{\textcolor{white}{\textbf{\endnotemark}}}}



\endnotetext{\vspace{-1.75\baselineskip}
\begin{tcolorbox}[blanker, breakable, enhanced, before skip = 5mm, after skip = 5mm, colframe=\currentcolor,left=3mm, right=3mm, top=6mm,segmentation style={draw=none,fill=none}, bottom=3mm, colback=white, %colframe=box1,
borderline west={3mm}{0pt}{white}, toprule=1pt, bottomrule=1pt, rightrule=1pt, leftrule=-1pt, outer arc = 3mm, arc = 3mm,
fonttitle=\bfseries\large,  % Should be 12 pt. Don't know how to set this.
title= Objetivos Específicos,  
boxed title style={empty,arc=0pt,outer arc=0pt,boxrule=0pt}, 
attach boxed title to top left={xshift=3mm , yshift=.5mm-.8pt},
underlay boxed title={
\draw[line width=3pt,white] ($(frame.north west)+(0,-1pt)$)--($(frame.north west)+(15pt,-1pt)$);
\filldraw [\currentcolor] ($(title.north west) + (0 mm,-3.50mm)$) -- ++(0.00mm,2.5mm) arc [start angle=180, end angle=90, y radius=2.00mm, x radius=2.00mm] -| ($(title.north east)+(0,-5mm)$) arc [start angle=0, end angle=-90, x radius=2.00mm, y radius=2.00mm] -| cycle;
}, 
before upper={\setlength{\parskip}{6pt}\vspace{-15pt}\tcbsubtitle{#1}
\normalfont\setlist[enumerate]{label=OE\arabic*,wide}\setlist[itemize,1]{wide, label=\textcolor{\currentcolor}{\small$\blacksquare$}, topsep=5pt}},
% {\boxtitlefont\setlist[enumerate]{label=OE\arabic*,wide}\setlist[itemize,1]{wide, label=\textcolor{penaqua}{\small$\blacksquare$}}
% {\vspace{1em}\large\textbf{Levar o aluno a}}\normalsize},
    subtitle style = {
      frame empty,
      fontupper = \bfseries\large,
      colback = white,
      coltext = \currentcolor,
      arc = 8pt,
      top=0pt,
      bottom=0pt
    },
    middle=5pt,
    break at=-\baselineskip/0pt] #2 \end{tcolorbox}

}

\else

\marginpar{\notebox{\textcolor{white}{\textbf{\endnotemark}}}}


\endnotetext{\adjustbox{valign=t}
{\begin{tcolorbox}[blanker, breakable, enhanced, before skip = 5mm, after skip = 5mm, colframe=\currentcolor,left=3mm, right=3mm, top=6mm,segmentation style={draw=none,fill=none}, bottom=3mm, colback=white, %colframe=box1,
borderline west={3mm}{0pt}{white}, toprule=1pt, bottomrule=1pt, rightrule=1pt, leftrule=-1pt, outer arc = 3mm, arc = 3mm,
fonttitle=\bfseries\large,  % Should be 12 pt. Don't know how to set this.
title= Objetivos Específicos,  
boxed title style={empty,arc=0pt,outer arc=0pt,boxrule=0pt}, 
attach boxed title to top left={xshift=3mm , yshift=.5mm-.8pt},
underlay boxed title={
\draw[line width=3pt,white] ($(frame.north west)+(0,-1pt)$)--($(frame.north west)+(15pt,-1pt)$);
\filldraw [\currentcolor] ($(title.north west) + (0 mm,-3.50mm)$) -- ++(0.00mm,2.5mm) arc [start angle=180, end angle=90, y radius=2.00mm, x radius=2.00mm] -| ($(title.north east)+(0,-5mm)$) arc [start angle=0, end angle=-90, x radius=2.00mm, y radius=2.00mm] -| cycle;
}, 
before upper={\setlength{\parskip}{6pt}\vspace{-15pt}\tcbsubtitle{#1}
\normalfont\setlist[enumerate]{label=OE\arabic*,wide}\setlist[itemize,1]{wide, label=\textcolor{\currentcolor}{\small$\blacksquare$}, topsep=5pt}},
% {\boxtitlefont\setlist[enumerate]{label=OE\arabic*,wide}\setlist[itemize,1]{wide, label=\textcolor{penaqua}{\small$\blacksquare$}}
% {\vspace{1em}\large\textbf{Levar o aluno a}}\normalsize},
    subtitle style = {
      frame empty,
      fontupper = \bfseries\large,
      colback = white,
      coltext = \currentcolor,
      arc = 8pt,
      top=0pt,
      bottom=0pt
    },
    middle=5pt,
    break at=-\baselineskip/0pt] #2 \end{tcolorbox}}

}
\fi}
\fi}
{}



%%%%%%%%%%%%% CAIXA DE ORIENTAÇÃO


\newenvironment{sugestions}[4]
{\ifnum #3=1   
{\ifnum #4=1

\marginpar{\vspace{.05em}\begin{tcolorbox}[blanker, enhanced, enforce breakable, before skip = 5mm, after skip = 5mm, colframe=\currentcolor, left=3mm, top=6mm,segmentation style={draw=none,fill=none}, bottom=3mm, colback=white, %colframe=box1,
borderline east={2mm}{0pt}{white}, toprule=1pt, bottomrule=1pt, leftrule=1pt, rightrule=-1pt, outer arc = 3mm, arc = 3mm,
fonttitle=\bfseries\large,  % Should be 12 pt. Don't know how to set this.
title= Sugestões e discussões, 
boxed title style={empty,arc=0pt,outer arc=0pt,boxrule=0pt}, 
attach boxed title to top right={ , yshift=.5mm-.2pt},
underlay boxed title={
\draw[line width=3pt,white] ($(frame.north east)+(0,-1pt)$)--($(frame.north east)+(-15pt,-1pt)$);
\filldraw [\currentcolor] ($(title.north west) + (0 mm,-3.50mm)$) -- ++(0.00mm,2.5mm) arc [start angle=180, end angle=90, y radius=2.00mm, x radius=2.00mm] -| ($(title.north east)+(0,-5mm)$) arc [start angle=0, end angle=-90, x radius=2.00mm, y radius=2.00mm] -| cycle;
}, 
before upper={\setlength{\parskip}{6pt}\vspace{-15pt}\tcbsubtitle{#1}
\normalfont\setlist[enumerate]{wide}\setlist[itemize,1]{wide, label=\textcolor{\currentcolor}{\small$\blacksquare$}, topsep=5pt}},
% {\boxtitlefont\setlist[enumerate]{label=OE\arabic*,wide}\setlist[itemize,1]{wide, label=\textcolor{penaqua}{\small$\blacksquare$}}
% {\vspace{1em}\large\textbf{Levar o aluno a}}\normalsize},
    subtitle style = {
      frame empty,
      fontupper = \bfseries\large,
      colback = white,
      coltext = \currentcolor,
      arc = 8pt,
      top=0pt,
      bottom=0pt
    },
    middle=5pt,
    break at=-\baselineskip/0pt] #2 \end{tcolorbox}}
    

\else


\marginpar{\vspace{.05em}\begin{tcolorbox}[blanker, enhanced, enforce breakable, before skip = 5mm, after skip = 5mm, colframe=\currentcolor, right=3mm, top=6mm,segmentation style={draw=none,fill=none}, bottom=3mm, colback=white, %colframe=box1,
borderline west={2mm}{0pt}{white}, toprule=1pt, bottomrule=1pt, rightrule=1pt, leftrule=-1pt, outer arc = 3mm, arc = 3mm,
fonttitle=\bfseries\large,  % Should be 12 pt. Don't know how to set this.
title= Sugestões e discussões, 
boxed title style={empty,arc=0pt,outer arc=0pt,boxrule=0pt}, 
attach boxed title to top left={ , yshift=.5mm-.2pt},
underlay boxed title={
\draw[line width=3pt,white] ($(frame.north west)+(0,-1pt)$)--($(frame.north west)+(15pt,-1pt)$);
\filldraw [\currentcolor] ($(title.north west) + (0 mm,-3.50mm)$) -- ++(0.00mm,2.5mm) arc [start angle=180, end angle=90, y radius=2.00mm, x radius=2.00mm] -| ($(title.north east)+(0,-5mm)$) arc [start angle=0, end angle=-90, x radius=2.00mm, y radius=2.00mm] -| cycle;
}, 
before upper={\setlength{\parskip}{6pt}\vspace{-15pt}\tcbsubtitle{#1}
\normalfont\setlist[enumerate]{wide}\setlist[itemize,1]{wide, label=\textcolor{\currentcolor}{\small$\blacksquare$}, topsep=5pt}},
% {\boxtitlefont\setlist[enumerate]{label=OE\arabic*,wide}\setlist[itemize,1]{wide, label=\textcolor{penaqua}{\small$\blacksquare$}}
% {\vspace{1em}\large\textbf{Levar o aluno a}}\normalsize},
    subtitle style = {
      frame empty,
      fontupper = \bfseries\large,
      colback = white,
      coltext = \currentcolor,
      arc = 8pt,
      top=0pt,
      bottom=0pt
    },
    middle=5pt,
    break at=-\baselineskip/0pt] #2 \end{tcolorbox}}
\fi}


\else
{\ifnum #4=9

\marginpar{\notebox{\textcolor{white}{\textbf{\endnotemark}}}\vspace{.05em}}

\endnotetext{
{\vspace{-2\baselineskip}
\begin{tcolorbox}[blanker, enhanced, breakable, before skip = 5mm, after skip = 5mm, colframe=\currentcolor,left=3mm, right=3mm, top=6mm,segmentation style={draw=none,fill=none}, bottom=3mm, colback=white, %colframe=box1,
borderline west={3mm}{0pt}{white}, toprule=1pt, bottomrule=1pt, rightrule=1pt, leftrule=-1pt, outer arc = 3mm, arc = 3mm,
fonttitle=\bfseries\large,  % Should be 12 pt. Don't know how to set this.
title= Sugestões e discussões,  
boxed title style={empty,arc=0pt,outer arc=0pt,boxrule=0pt}, 
attach boxed title to top left={xshift=3mm , yshift=.5mm-.2pt},
underlay boxed title={
\draw[line width=3pt,white] ($(frame.north west)+(0,-1pt)$)--($(frame.north west)+(15pt,-1pt)$);
\filldraw [\currentcolor] ($(title.north west) + (0 mm,-3.50mm)$) -- ++(0.00mm,2.5mm) arc [start angle=180, end angle=90, y radius=2.00mm, x radius=2.00mm] -| ($(title.north east)+(0,-5mm)$) arc [start angle=0, end angle=-90, x radius=2.00mm, y radius=2.00mm] -| cycle;
}, 
before upper={\setlength{\parskip}{6pt}\vspace{-15pt}\tcbsubtitle{#1}
\normalfont\setlist[enumerate]{wide}\setlist[itemize,1]{wide, label=\textcolor{\currentcolor}{\small$\blacksquare$}, topsep=5pt}},
% {\boxtitlefont\setlist[enumerate]{label=OE\arabic*,wide}\setlist[itemize,1]{wide, label=\textcolor{penaqua}{\small$\blacksquare$}}
% {\vspace{1em}\large\textbf{Levar o aluno a}}\normalsize},
    subtitle style = {
      frame empty,
      fontupper = \bfseries\large,
      colback = white,
      coltext = \currentcolor,
      arc = 8pt,
      top=0pt,
      bottom=0pt
    },
    middle=5pt,
    break at=-\baselineskip/0pt] #2 \end{tcolorbox}}


}

\else

\marginpar{\notebox{\textcolor{white}{\textbf{\endnotemark}}}}

\endnotetext{\adjustbox{valign=t}
{
\begin{tcolorbox}[blanker, enhanced, breakable, before skip = 5mm, after skip = 5mm, colframe=\currentcolor,left=3mm, right=3mm, top=6mm,segmentation style={draw=none,fill=none}, bottom=3mm, colback=white, %colframe=box1,
borderline west={3mm}{0pt}{white}, toprule=1pt, bottomrule=1pt, rightrule=1pt, leftrule=-1pt, outer arc = 3mm, arc = 3mm,
fonttitle=\bfseries\large,  % Should be 12 pt. Don't know how to set this.
title= Sugestões e discussões,  
boxed title style={empty,arc=0pt,outer arc=0pt,boxrule=0pt}, 
attach boxed title to top left={xshift=3mm , yshift=.5mm-.2pt},
underlay boxed title={
\draw[line width=3pt,white] ($(frame.north west)+(0,-1pt)$)--($(frame.north west)+(15pt,-1pt)$);
\filldraw [\currentcolor] ($(title.north west) + (0 mm,-3.50mm)$) -- ++(0.00mm,2.5mm) arc [start angle=180, end angle=90, y radius=2.00mm, x radius=2.00mm] -| ($(title.north east)+(0,-5mm)$) arc [start angle=0, end angle=-90, x radius=2.00mm, y radius=2.00mm] -| cycle;
}, 
before upper={\setlength{\parskip}{6pt}\vspace{-15pt}\tcbsubtitle{#1}
\normalfont\setlist[enumerate]{wide}\setlist[itemize,1]{wide, label=\textcolor{\currentcolor}{\small$\blacksquare$}, topsep=5pt}},
% {\boxtitlefont\setlist[enumerate]{label=OE\arabic*,wide}\setlist[itemize,1]{wide, label=\textcolor{penaqua}{\small$\blacksquare$}}
% {\vspace{1em}\large\textbf{Levar o aluno a}}\normalsize},
    subtitle style = {
      frame empty,
      fontupper = \bfseries\large,
      colback = white,
      coltext = \currentcolor,
      arc = 8pt,
      top=0pt,
      bottom=0pt
    },
    middle=5pt,
    break at=-\baselineskip/0pt] #2 \end{tcolorbox}}
}
\fi}
\fi}


\setlist[enumerate]{label=\protect\bfseries\textcolor{\currentcolor}{\alph*)}}



%%%%%%%%%%%%% CAIXA DE RESPOSTA
\newenvironment{answer}[3]
{\ifnum #3=1   
\marginpar{\begin{tcolorbox}[enhanced, breakable,empty, coltitle=penazul, colback=white!50, left=0mm, right=0mm, boxrule=0mm, boxsep=0mm, colframe=penaqua,colbacktitle=white, before title ={\boxtitlefont \textcolor{\currentcolor}{Solução: }}, title= \textcolor{grafite}{#1},
before upper={ \setlength{\parskip}{6pt}\setlist[enumerate]{label=wide,left=0pt,label=\protect\bfseries\textcolor{\currentcolor}{\alph*)}}},
    subtitle style = {
      frame empty,
      fontupper = \boxtitlefont,
      colback = white,
      coltext = grafite,
      top=1em,
      before upper = \boxtitlefont \textcolor{\currentcolor}{Solução: }}]
#2\end{tcolorbox}}
\else
{\ifnum #3=9

\marginpar{\notebox{\textcolor{white}{\textbf{\endnotemark}}}}

\endnotetext{\vspace{-1.35\baselineskip}
\begin{tcolorbox}[enhanced, breakable,empty, coltitle=penazul, colback=white, left=0mm, right=0mm, boxrule=0mm, boxsep=0mm, colframe=penaqua,colbacktitle=white, before title ={\boxtitlefont \textcolor{\currentcolor}{Solução: }}, title= \textcolor{grafite}{#1}, before upper={\setlength{\parskip}{6pt}\vspace{.3em}
\setlist[enumerate]{label=wide,left=0pt,label=\protect\bfseries\textcolor{\currentcolor}{\alph*)}}
\setlist[enumerate,2]{label=\protect\bfseries\textcolor{\currentcolor}{\roman*)}}
},
    subtitle style = {
      frame empty,
      fontupper = \boxtitlefont,
      colback = white,
      coltext = grafite,
      top=1em,
      before upper = \boxtitlefont \textcolor{\currentcolor}{Solução: }}]
#2\end{tcolorbox}}

\else
\marginpar{\notebox{\textcolor{white}{\textbf{\endnotemark}}}}

\endnotetext{\adjustbox{valign=t}
{

\begin{tcolorbox}[enhanced, breakable,empty, coltitle=penazul, colback=white, left=0mm, right=0mm, boxrule=0mm, boxsep=0mm, colframe=penaqua,colbacktitle=white, before title ={\boxtitlefont \textcolor{\currentcolor}{Solução: }}, title= \textcolor{grafite}{#1}, before upper={\setlength{\parskip}{6pt}\vspace{.3em}
\setlist[enumerate]{label=wide,left=0pt,label=\protect\bfseries\textcolor{\currentcolor}{\alph*)}}
\setlist[enumerate,2]{label=\protect\bfseries\textcolor{\currentcolor}{\roman*)}}
},
    subtitle style = {
      frame empty,
      fontupper = \boxtitlefont,
      colback = white,
      coltext = grafite,
      top=1em,
      before upper = \boxtitlefont \textcolor{\currentcolor}{Solução: }}]
#2\end{tcolorbox}}
}

\fi}
\fi
}


\newcommand\notas[1]{
	\notebox{\textcolor{white}{\textbf{\endnotemark}}}
	\endnotetext
	{#1}}
\newcommand\mnotas[1]{
	\marginpar{
	\notebox{\textcolor{white}{\textbf{\endnotemark}}}
	\endnotetext
	{#1}}}


\newcommand\notasfinais{	
	\newgeometry{textwidth=840pt,outer=50pt, top=2cm}
	\pagestyle{plain}
	\def\currentcolor{cor1}
	\section{\huge Notas}
	\setlength{\columnsep}{1cm}
	\begin{multicols}{2}\raggedcolumns
	\theendnotes
	\clearpage
	\end{multicols}
	\restoregeometry
}

%%%%%%%%%%%%% CAIXA DE HABILIDADES
\newtcolorbox{habilities}[1]{ breakable,colback=penaqua!20!white,colbacktitle=penaqua!20!white,
colframe=penaqua!20!white, left =2mm, right=2mm, boxrule=0pt, top=1mm, arc=3mm, toptitle=2mm, coltitle=black, fonttitle=\bfseries, arc is angular = northwest, sharp corners = northeast, sharp corners = southwest, arc is angular = southeast, subtitle style={coltitle=black,fonttitle=\bfseries},
title=#1}

\newenvironment{OES}{\begin{enumerate}[label=\titem{OE \arabic*} -, topsep=2.5pt,itemsep=0pt]}{\end{enumerate}}

%%%%%%%%%%%%%%% NOTA DE RODAPÉ
\newcommand{\notarodape}[1]{\footnotemark\footnotetext{\belowbaseline[-\ht\strutbox]{\begin{tcolorbox}[colback=grafite!90!white,colbacktitle=grafite!90!white,
colframe=grafite!90!white,colupper=white, left=0pt, right=0pt,top=-2pt,bottom=-2pt, hbox, sharp corners = northeast, sharp corners = southwest]
{\textbf{#1}}\end{tcolorbox}}}}

\geometry{reset, paperwidth=950pt,paperheight=782pt,textwidth=450pt,
marginparsep=30pt,headsep=19pt,top=45pt,inner=50pt,outer=440pt,
marginparwidth=350pt,textheight=686pt}





\endinput

